% Compile this document using XeLaTeX or LuaLaTeX
\documentclass[12pt,a4paper,oneside]{report}

% ─── Packages ─────────────────────────────────────────────────────────────────
\usepackage{fontspec}
\usepackage{polyglossia}
\setdefaultlanguage{english}
\setotherlanguage{hindi}
\newfontfamily\hindifont[Script=Devanagari]{Noto Serif Devanagari} % Change to Mangal, Arial Unicode MS, or FreeSerif if Noto is not installed

\usepackage{amsmath,amssymb,amsfonts,amsthm}
\usepackage{graphicx}
\usepackage{geometry}
\usepackage{setspace}
\usepackage{fancyhdr}
\usepackage{titlesec}
\usepackage{tocloft}
\usepackage{booktabs}
\usepackage{multirow}
\usepackage{array}
\usepackage{longtable}
\usepackage{tabularx}
\usepackage{url}
\usepackage{hyperref}
\usepackage{xcolor}
\usepackage{caption}
\usepackage{algorithm}
\usepackage{algpseudocode}
\usepackage{listings}
\usepackage{tikz}
\usetikzlibrary{shapes.geometric,arrows.meta,positioning,calc}
\usepackage{enumitem}
\usepackage{float}

% ─── Page geometry ────────────────────────────────────────────────────────────
\geometry{a4paper,top=1in,bottom=1in,left=1.5in,right=1in,headheight=15pt}

% ─── Spacing ──────────────────────────────────────────────────────────────────
\onehalfspacing

% ─── Hyperref ─────────────────────────────────────────────────────────────────
\hypersetup{colorlinks=true,linkcolor=black,citecolor=black,
  urlcolor=blue!60!black,
  pdftitle={ZK-Auth Thesis},pdfauthor={Abhay Dandge}}

% ─── Header/Footer ────────────────────────────────────────────────────────────
\pagestyle{fancy}
\fancyhf{}
\fancyfoot[C]{\small\thepage}
\renewcommand{\headrulewidth}{0pt}
\fancypagestyle{plain}{
  \fancyhf{}
  \fancyfoot[C]{\small\thepage}
  \renewcommand{\headrulewidth}{0pt}
}

% ─── TOC Dots for Chapters ────────────────────────────────────────────────────
\renewcommand{\cftchapleader}{\cftdotfill{\cftdotsep}}

% ─── Initial Chapter/section format ───────────────────────────────────────────
\titleformat{\chapter}[display]
  {\normalfont\Large\bfseries\raggedleft}{CHAPTER~\thechapter}{12pt}{\large\MakeUppercase}
\titlespacing*{\chapter}{0pt}{-10pt}{30pt}

\titleformat{name=\chapter,numberless}[block]
  {\normalfont\Large\bfseries\raggedleft}{}{0pt}{\large\MakeUppercase}

\titleformat{\section}{\normalfont\large\bfseries}{\thesection}{1em}{}
\titleformat{\subsection}{\normalfont\normalsize\bfseries}{\thesubsection}{1em}{}

% ─── TikZ styles ──────────────────────────────────────────────────────────────
\tikzset{
  block/.style={rectangle,draw=black!70,fill=blue!8,rounded corners=3pt,
    minimum height=0.8cm,align=center,font=\scriptsize},
  actor/.style={rectangle,draw=black!80,fill=gray!10,rounded corners=4pt,
    minimum width=2.2cm,minimum height=0.8cm,align=center,font=\scriptsize\bfseries},
  procbox/.style={rectangle,draw=black!60,fill=green!8,rounded corners=2pt,
    minimum width=2.5cm,minimum height=0.7cm,align=center,font=\scriptsize},
  barlabel/.style={font=\scriptsize},
}

% Simple manual bar-chart macro
\newcommand{\simplebar}[6]{%
  \draw[fill=#5,draw=black!60] (#1,0) rectangle ++(#3,#4);%
  \node[barlabel,above] at (#1+0.5*#3,#4) {#6};%
}
\newcommand{\simplebarneg}[5]{%
  \draw[fill=#4,draw=black!60] (#1,0) rectangle ++(#3,#2);%
  \node[barlabel,below] at (#1+0.5*#3,#2) {#5};%
}

% ─── Theorem environments ─────────────────────────────────────────────────────
\newtheorem{definition}{Definition}[chapter]
\newtheorem{theorem}{Theorem}[chapter]
\newtheorem{property}{Property}[chapter]
\newtheorem{lemma}{Lemma}[chapter]
\newtheorem{corollary}{Corollary}[chapter]

% ─── Listings ─────────────────────────────────────────────────────────────────
\lstset{basicstyle=\ttfamily\small,breaklines=true,frame=single,
  numbers=left,numberstyle=\tiny,
  keywordstyle=\bfseries\color{blue!60!black},
  commentstyle=\itshape\color{green!40!black}}

%==============================================================================
\begin{document}
%==============================================================================

%------------------------------------------------------------------------------
% TITLE PAGE
%------------------------------------------------------------------------------
\begin{titlepage}
\vspace*{\fill}
\begin{center}
{\Large\bfseries
ZERO-KNOWLEDGE PROOF-BASED AUTHENTICATION FRAMEWORK\\[4pt]
WITH SELECTIVE DISCLOSURE AND BEHAVIORAL BIOMETRICS\\[4pt]
FOR CLOUD COMPUTING
}\\[0.6cm]
\vfill
{\large A}\\[0.2cm]
{\Large\bfseries THESIS}\\[0.3cm]
{\normalsize\itshape
Submitted in partial fulfillment of the requirements for the award of degree of}\\[0.3cm]
{\large\bfseries MASTERS OF TECHNOLOGY}\\[0.2cm]
{\normalsize\itshape In}\\[0.2cm]
{\large\bfseries ARTIFICIAL INTELLIGENCE}\\[0.4cm]
\vfill
{\normalsize\itshape Submitted by}\\[0.2cm]
{\large\bfseries ABHAY DANDGE}\\[0.2cm]
{\normalsize Scholar Number: \textbf{24215011117}}\\[0.4cm]
\vfill
{\normalsize\itshape Under the Guidance of}\\[0.2cm]
{\normalsize\bfseries Dr.~Namita Tiwari}\\
{\normalsize Assistant Professor, Department of CSE, MANIT Bhopal}\\[0.2cm]
{\normalsize\bfseries Dr.~Meenu Chawla}\\
{\normalsize Professor, Department of CSE, MANIT Bhopal}\\[0.5cm]
\vfill
\includegraphics[width=2.5cm]{manit_logo.png}\\[0.2cm]
{\normalsize\bfseries Centre of Artificial Intelligence}\\[0.2cm]
{\large\bfseries MAULANA AZAD NATIONAL INSTITUTE OF TECHNOLOGY BHOPAL}\\[0.1cm]
{\normalsize\bfseries JULY 2026}
\end{center}
\vspace*{\fill}
\end{titlepage}

%------------------------------------------------------------------------------
% FRONT MATTER
%------------------------------------------------------------------------------
\clearpage
\pagenumbering{roman}
\setcounter{page}{2}

% DECLARATION
\clearpage
\phantomsection
\addcontentsline{toc}{chapter}{Declaration}
\begin{center}
{\large\bfseries MAULANA AZAD NATIONAL INSTITUTE OF TECHNOLOGY BHOPAL}\\[0.1cm]
{\normalsize (An Institute of National Importance)}\\[0.1cm]
{\normalsize Bhopal - 462003}\\[0.4cm]
\includegraphics[width=2.5cm]{manit_logo.png}\\[0.4cm]
{\normalsize\bfseries Centre of Artificial Intelligence}\\[0.6cm]
{\Large\bfseries\underline{Declaration}}
\end{center}

\vspace{0.5cm}
\noindent I hereby declare that the work presented in the dissertation entitled
\textbf{``Zero-Knowledge Proof-Based Authentication Framework with Selective
Disclosure and Behavioral Biometrics for Cloud Computing''} is submitted in
partial fulfillment of the requirements for the award of the degree of
\textbf{Masters of Technology} in \textbf{Artificial Intelligence}.
It is an authentic documentation of my own original work carried out from
July~2024 to July~2026 under the guidance of \textbf{Dr.~Namita Tiwari},
Assistant Professor, and co-supervisor \textbf{Dr.~Meenu Chawla}, Professor,
Department of CSE, MANIT Bhopal.

Wherever contributions of others are involved, every effort is made to indicate
this clearly with due reference to the literature. The matter embodied in this
dissertation has not been presented or submitted for any other degree.

\vspace{2.5cm}
\begin{flushright}
\rule{6cm}{0.4pt}\\[5pt]
\textbf{Abhay Dandge}\\
Scholar Number: 24215011117\\
M.Tech (Artificial Intelligence)\\
MANIT, Bhopal
\end{flushright}

% CERTIFICATE
\clearpage
\phantomsection
\addcontentsline{toc}{chapter}{Certificate}
\begin{center}
{\large\bfseries MAULANA AZAD NATIONAL INSTITUTE OF TECHNOLOGY BHOPAL}\\[0.1cm]
{\normalsize (An Institute of National Importance)}\\[0.1cm]
{\normalsize Bhopal - 462003}\\[0.4cm]
\includegraphics[width=2.5cm]{manit_logo.png}\\[0.4cm]
{\normalsize\bfseries Centre of Artificial Intelligence}\\[0.6cm]
{\Large\bfseries\underline{Certificate}}
\end{center}

\vspace{0.5cm}
\noindent This is to certify that the dissertation entitled \textbf{``Zero-Knowledge
Proof-Based Authentication Framework with Selective Disclosure and Behavioral
Biometrics for Cloud Computing''} submitted by \textbf{Abhay Dandge},
Scholar No.~24215011117, in partial fulfillment of the requirements for the
degree of \textbf{Masters of Technology} in Artificial Intelligence, is an
authentic work carried out under our supervision and guidance from
July~2024 to July~2026.

\vspace{3cm}
\noindent
\begin{minipage}[t]{0.45\textwidth}
\rule{5.5cm}{0.4pt}\\[3pt]
\textbf{Dr.~Namita Tiwari}\\
Supervisor\\
Assistant Professor, Dept. of CSE\\
MANIT, Bhopal
\end{minipage}
\hfill
\begin{minipage}[t]{0.45\textwidth}
\raggedleft
\rule{5.5cm}{0.4pt}\\[3pt]
\textbf{Dr.~Meenu Chawla}\\
Co-supervisor\\
Professor, Dept. of CSE\\
MANIT, Bhopal
\end{minipage}


% ─── Redefine Chapter Formats for Line Under Headings ─────────────────────────
\titleformat{name=\chapter,numberless}[block]
  {\normalfont\Large\bfseries\raggedleft}{}{0pt}{\large\MakeUppercase}[\vspace{1ex}\titlerule]

\titleformat{\chapter}[display]
  {\normalfont\Large\bfseries\raggedleft}{CHAPTER~\thechapter}{12pt}{\large\MakeUppercase}[\vspace{1ex}\titlerule]
% ──────────────────────────────────────────────────────────────────────────────


% ACKNOWLEDGEMENT
\clearpage
\chapter*{Acknowledgement}
\addcontentsline{toc}{chapter}{Acknowledgement}

I would like to begin by expressing my profound sense of gratitude towards my
supervisor \textbf{Dr.~Namita Tiwari}, Assistant Professor, and co-supervisor
\textbf{Dr.~Meenu Chawla}, Professor, Department of CSE, MANIT Bhopal, for
their valuable guidance, support, and encouragement throughout this research.
Their readiness for consultation at all times, their educative comments, and
their constant concern for quality have been invaluable.

I express my sincere gratitude to our Director, \textbf{Dr.~K.K.~Shukla}, for
providing ample resources and a conducive research environment. I extend my
thanks to \textbf{Dr.~R.K.~Pateriya}, Head of Department of CSE, for
administrative support and for providing access to departmental laboratories.

I am grateful to my co-author \textbf{Sameeksha Prasad} for collaborative
research discussions that contributed to the published works underpinning this
thesis. I also thank my fellow research scholars and classmates for stimulating
academic discussions and collegial support throughout the programme.

Finally, I am deeply indebted to my parents and family for everything they have
given me. They have always stood by me with constant support, encouragement, and
love through every phase of this research.

\vspace{2.5cm}
\noindent
\begin{minipage}[t]{0.45\textwidth}
\textbf{Date:}\\[0.4cm]
\textbf{Place:} Bhopal
\end{minipage}
\hfill
\begin{minipage}[t]{0.45\textwidth}
\raggedleft
\textbf{Abhay Dandge}\\
Scholar Number: 24215011117
\end{minipage}

% ABSTRACT (ENGLISH)
\clearpage
\chapter*{Abstract}
\addcontentsline{toc}{chapter}{Abstract}

Authentication security in cloud environments remains a critical vulnerability. Traditional mechanisms---passwords, multi-factor authentication (MFA), and public key infrastructure (PKI)---rely on a \textit{verify-by-reveal} paradigm. Sensitive credentials must be transmitted or stored by the verifier, creating massive attack surfaces highly exploitable through database breaches and credential stuffing.

This thesis introduces \textsc{ZK-Auth}, a cryptographic infrastructure resolving the verify-by-reveal problem. \textsc{ZK-Auth} operates as an OAuth 2.0 Authorization Server where the standard password authentication step is replaced by a zero-knowledge proof. Three core mechanisms realize this design. \textit{First}, users prove knowledge of a registration secret strictly on the client side using Groth16 zk-SNARKs over the BN254 elliptic curve, compiled to WebAssembly. The resulting proof authorizes a short-lived OAuth authorization code. Poseidon hash-based nullifiers simultaneously prevent replay attacks.

\textit{Second}, selective credential disclosure is achieved via depth-8 Poseidon Merkle trees integrated with W3C Verifiable Credentials. Holders prove individual attribute predicates (e.g., age $\geq$ 18) without exposing raw attribute values. \textit{Third}, \textsc{ZK-Auth} incorporates a two-layer LSTM behavioral biometric model providing background session-risk scoring derived from user telemetry, triggering step-up re-authentication upon detecting anomalous behavior.

The architecture supports multi-tenant institutional issuance via Ed25519 signing keypairs. \textsc{ZK-Auth} is implemented end-to-end and evaluated across Node.js, Chrome, and two Android devices. It achieves a median end-to-end authentication latency of 132.5\,ms---23\% faster than an Argon2id password-hashing baseline---with a 723-byte proof. The behavioral biometric subsystem, evaluated on the public Balabit Mouse Dynamics dataset, achieves an AUC-ROC of 0.816 with a 6.15\% false-acceptance rate. This zero-PII architecture structurally complies with India's DPDP Act 2023 and the GDPR.

% ABSTRACT (HINDI)
\clearpage
\chapter*{\texthindi{संक्षेप}}
\addcontentsline{toc}{chapter}{Abstract (Hindi)}

\begin{hindi}
क्लाउड परिवेशों में प्रमाणीकरण सुरक्षा अब भी एक गंभीर भेद्यता बनी हुई है। पारंपरिक तंत्र—पासवर्ड, बहु-कारक प्रमाणीकरण (\textenglish{MFA}), और सार्वजनिक-कुंजी अवसंरचना (\textenglish{PKI})—एक ``\textenglish{verify-by-reveal}'' प्रतिमान पर निर्भर करते हैं। संवेदनशील क्रेडेंशियल्स को सत्यापनकर्ता द्वारा प्रेषित या संग्रहीत करना आवश्यक होता है, जिससे विशाल आक्रमण सतहें बनती हैं, जिनका डेटाबेस उल्लंघनों और क्रेडेंशियल स्टफिंग के माध्यम से अत्यधिक शोषण किया जा सकता है।

यह शोधप्रबंध \textenglish{\textsc{ZK-Auth}} प्रस्तुत करता है, जो एक ऐसी क्रिप्टोग्राफिक अवसंरचना है जो ``\textenglish{verify-by-reveal}'' की समस्या का समाधान करती है। \textenglish{\textsc{ZK-Auth}} एक \textenglish{OAuth 2.0 Authorization Server} के रूप में कार्य करता है, जहाँ मानक पासवर्ड प्रमाणीकरण चरण को शून्य-ज्ञान प्रमाण (\textenglish{zero-knowledge proof}) द्वारा प्रतिस्थापित किया जाता है। इस डिज़ाइन को तीन मुख्य तंत्र साकार करते हैं। पहला, उपयोगकर्ता \textenglish{Groth16 zk-SNARKs} का उपयोग करते हुए, \textenglish{BN254} एलिप्टिक वक्र पर, \textenglish{WebAssembly} में संकलित प्रणाली के माध्यम से, केवल क्लाइंट-पक्ष पर पंजीकरण रहस्य (\textenglish{registration secret}) के ज्ञान का प्रमाण प्रस्तुत करते हैं। परिणामी प्रमाण एक अल्पकालिक \textenglish{OAuth} प्राधिकरण कोड को अधिकृत करता है। \textenglish{Poseidon} हैश-आधारित नलिफ़ायर्स (\textenglish{nullifiers}) एक साथ रिप्ले आक्रमणों को रोकते हैं।

दूसरा, चयनात्मक क्रेडेंशियल प्रकटीकरण (\textenglish{selective credential disclosure}) गहराई-8 \textenglish{Poseidon Merkle} वृक्षों के माध्यम से प्राप्त किया जाता है, जिन्हें \textenglish{W3C Verifiable Credentials} के साथ एकीकृत किया गया है। धारक कच्चे गुण मानों (\textenglish{raw attribute values}) को उजागर किए बिना व्यक्तिगत गुण-आधारित विधेयों (जैसे, आयु $\geq$ 18) का प्रमाण प्रस्तुत करते हैं। तीसरा, \textenglish{\textsc{ZK-Auth}} एक द्वि-स्तरीय \textenglish{LSTM} व्यवहारिक बायोमेट्रिक मॉडल को सम्मिलित करता है, जो उपयोगकर्ता टेलीमेट्री से प्राप्त पृष्ठभूमि सत्र-जोखिम स्कोरिंग प्रदान करता है और असामान्य व्यवहार का पता चलने पर उन्नत पुनः-प्रमाणीकरण (\textenglish{step-up re-authentication}) को सक्रिय करता है।

यह वास्तुकला \textenglish{Ed25519} हस्ताक्षर कुंजी-युग्मों (\textenglish{signing keypairs}) के माध्यम से बहु-किरायेदार (\textenglish{multi-tenant}) संस्थागत निर्गमन का समर्थन करती है। \textenglish{\textsc{ZK-Auth}} को पूर्णतः कार्यान्वित किया गया है और \textenglish{Node.js}, \textenglish{Chrome} तथा दो \textenglish{Android} उपकरणों पर मूल्यांकित किया गया है। यह 132.5 मिलीसेकंड की माध्य अंत-से-अंत प्रमाणीकरण विलंबता प्राप्त करता है—जो \textenglish{Argon2id} पासवर्ड-हैशिंग आधाररेखा की तुलना में 23\% तेज़ है—और इसका प्रमाण आकार 723 बाइट है। व्यवहारिक बायोमेट्रिक उपप्रणाली, जिसे सार्वजनिक \textenglish{Balabit Mouse Dynamics} डेटासेट पर मूल्यांकित किया गया, 0.816 का \textenglish{AUC-ROC} तथा 6.15\% की गलत-स्वीकृति दर (\textenglish{false-acceptance rate}) प्राप्त करती है। यह शून्य-\textenglish{PII} वास्तुकला भारत के \textenglish{DPDP} अधिनियम 2023 तथा \textenglish{GDPR} के साथ संरचनात्मक रूप से अनुपालन करती है।
\end{hindi}

% TABLE OF CONTENTS / LISTS
%------------------------------------------------------------------------------
\clearpage
\tableofcontents

\clearpage
\phantomsection
\addcontentsline{toc}{chapter}{List of Tables}
\listoftables

\clearpage
\phantomsection
\addcontentsline{toc}{chapter}{List of Figures}
\listoffigures

% LIST OF ABBREVIATIONS
\clearpage
\chapter*{List of Abbreviations}
\addcontentsline{toc}{chapter}{List of Abbreviations}

\begin{longtable}{p{3.5cm} p{10cm}}
\toprule
\textbf{Abbreviation} & \textbf{Full Form} \\
\midrule
\endhead
ABE     & Attribute-Based Encryption \\
API     & Application Programming Interface \\
AUC     & Area Under the (ROC) Curve \\
BN254   & Barreto-Naehrig 254-bit elliptic curve \\
CA      & Certificate Authority \\
CP-ABE  & Ciphertext-Policy Attribute-Based Encryption \\
CRS     & Common Reference String \\
DID     & Decentralized Identifier \\
DPDP    & Digital Personal Data Protection (Act, India 2023) \\
EMA     & Exponential Moving Average \\
EUF-CMA & Existential Unforgeability under Chosen Message Attack \\
FAR     & False Acceptance Rate \\
FRR     & False Rejection Rate \\
GDPR    & General Data Protection Regulation \\
gRPC    & Google Remote Procedure Call \\
IAM     & Identity and Access Management \\
IdP     & Identity Provider \\
IoT     & Internet of Things \\
JWT     & JSON Web Token \\
KZG     & Kate-Zaverucha-Goldberg polynomial commitment \\
LSTM    & Long Short-Term Memory \\
MFA     & Multi-Factor Authentication \\
NIZK    & Non-Interactive Zero-Knowledge \\
OAuth   & Open Authorization \\
PKI     & Public Key Infrastructure \\
PII     & Personally Identifiable Information \\
PPT     & Probabilistic Polynomial Time \\
QAP     & Quadratic Arithmetic Program \\
R1CS    & Rank-1 Constraint System \\
SAML    & Security Assertion Markup Language \\
SISMEMBER & Redis Set-Is-Member command \\
SNARK   & Succinct Non-interactive ARgument of Knowledge \\
STARK   & Scalable Transparent ARgument of Knowledge \\
SSI     & Self-Sovereign Identity \\
TEE     & Trusted Execution Environment \\
TTL     & Time To Live \\
VC      & Verifiable Credential \\
VP      & Verifiable Presentation \\
W3C     & World Wide Web Consortium \\
WASM    & WebAssembly \\
WebView & Embedded browser engine within a native mobile app \\
ZKP     & Zero-Knowledge Proof \\
ZK-SNARK & Zero-Knowledge Succinct Non-Interactive Argument of Knowledge \\
ZK-STARK & Zero-Knowledge Scalable Transparent Argument of Knowledge \\
\bottomrule
\end{longtable}

%==============================================================================
% MAIN MATTER
%==============================================================================
\clearpage
\pagenumbering{arabic}
\setcounter{page}{1}

%==============================================================================
\chapter{Introduction}
%==============================================================================

\section{Background and Motivation}

Cloud computing has transformed digital service delivery, with the global cloud
market exceeding \$600 billion annually by 2024. However, the shift to cloud
infrastructure introduces a critical security challenge: authenticating users in
environments where the infrastructure is not under direct organizational control.

Traditional authentication mechanisms---passwords, multi-factor authentication
(MFA), and public key infrastructure (PKI)---were designed for environments where
the authenticating server is trusted. In cloud deployments, this assumption often
does not hold. The \textit{verify-by-reveal} paradigm requires credentials to be
transmitted to or stored by the authenticating party, creating attack surfaces that
cloud-scale breaches have repeatedly exploited. Password database breaches expose
user credentials at scale; man-in-the-middle attacks intercept credentials in
transit; credential stuffing leverages leaked databases across services.

Zero-Knowledge Proofs (ZKPs) offer a mathematically rigorous resolution to this
challenge. A ZKP allows a \textit{prover} to convince a \textit{verifier} that a
statement is true without revealing any information beyond the truth of the
statement~\cite{gmr89}. Applied to authentication, a user can prove knowledge of a
registration secret without ever transmitting it---eliminating the primary attack
surface of credential theft. The verifier receives a cryptographic proof that is
sound (cannot be forged without knowledge of the secret) and zero-knowledge
(reveals nothing about the secret).

Despite their theoretical appeal, ZKPs have seen limited adoption in practical
authentication systems due to four challenges identified and addressed in this
thesis.

\textbf{Challenge 1 --- Proof complexity and cross-environment viability.}
Proof generation is computationally intensive, requiring efficient circuit
design and client-side execution within user-acceptable latency. Critically,
this latency bound is rarely characterized across the realistic spread of
client hardware (desktop browsers, flagship smartphones, mid-range tablets)
that a real user population presents; most prior ZKP-authentication
literature reports a single-environment benchmark only.

\textbf{Challenge 2 --- Attribute-level privacy.} Authentication in real
deployments requires proving specific \textit{attributes} of identity with
selective revelation---a capability basic ZKP authentication does not
provide.

\textbf{Challenge 3 --- Session continuity.} One-time authentication proofs
cannot protect against session hijacking; continuous identity assurance
mechanisms are required throughout a session lifetime, and these mechanisms
must themselves be evaluated against authentic (not synthetic) behavioral
data to be credible.

\textbf{Challenge 4 --- Empirical and theoretical rigor.} Many published ZKP
authentication systems either omit quantitative benchmarking entirely or
provide only qualitative security claims without a constraint-level
reduction tying the specific circuit design to the underlying proof
system's guarantees.

This thesis presents \textsc{ZK-Auth}, a complete cryptographic authentication
infrastructure addressing all four challenges in an integrated, empirically
evaluated, and formally analyzed deployable system.

\section{Published Research Underpinning This Thesis}

This thesis synthesizes and substantially extends two peer-reviewed publications
produced during the M.Tech programme:

\textbf{Publication 1~\cite{dandge2025zkpcloud}:} Dandge, A., Prasad, S., Tiwari,
N., and Chawla, M. ``Zero-Knowledge Proof for Authentication in Cloud Computing:
A Review.'' \textit{Springer Nature, FrontSci 2025}. This paper established the
architectural taxonomy of ZKP deployment patterns for cloud authentication,
compared ZKP families (zk-SNARKs, zk-STARKs, Bulletproofs, Sigma protocols) on
cloud-specific performance parameters, and identified the performance-auditability
tradeoff between stand-alone and ledger-anchored deployment. Table~\ref{tab:zkp_families}
in Chapter~2 reproduces this paper's central comparative finding, which directly
motivated the Groth16 design choice in \textsc{ZK-Auth}.

\textbf{Publication 2~\cite{prasad2025abereview}:} Prasad, S., Dandge, A., Tiwari,
N., and Chawla, M. ``A Comprehensive Review of Zero-Knowledge Proofs in
Attribute-Based Encryption Frameworks.'' \textit{Accepted, ICNDIA 2026 (IEEE)}.
This paper analyzed ZKP integration with ABE schemes for fine-grained access
control, established performance bounds of polynomial commitment schemes (KZG),
and identified hybrid ZKP models and hardware acceleration as critical research
directions. The selective-disclosure design in Chapter~4 of this thesis draws
directly on this paper's comparative analysis of ZK-compatible versus non-ZK
disclosure mechanisms.

Both publications surveyed the literature and proposed architectural
recommendations; neither implemented or empirically evaluated a deployable
system. This thesis closes that gap: it is the first work in this research
programme to implement, deploy, and benchmark a complete ZKP-based
authentication system across multiple real execution environments.

\section{Problem Statement}

Formally, the problem this thesis addresses can be stated as follows. Given
a population of users $\mathcal{U}$ and a set of relying-party services
$\mathcal{V}$, design an authentication protocol $\Pi$ such that:

\begin{enumerate}
  \item For any user $u \in \mathcal{U}$ and verifier $v \in \mathcal{V}$,
        $u$ can prove possession of a valid credential to $v$ without
        transmitting any secret material that, if intercepted or stored by
        $v$, would allow $v$ (or an attacker compromising $v$) to
        impersonate $u$ to a different verifier $v' \in \mathcal{V}$.
  \item $u$ can selectively disclose a predicate over a subset of identity
        attributes (e.g., $\text{age}(u) \geq 18$) to $v$ without revealing
        the underlying attribute value or any other attribute.
  \item Once authenticated, $u$'s session can be continuously monitored for
        anomalous behavior indicative of session hijacking, with a bounded
        detection-to-remediation window, without requiring $u$ to
        re-authenticate on every request.
  \item The protocol $\Pi$ must be empirically demonstrated to operate
        within human-acceptable latency bounds (conventionally,
        sub-second response time for an authentication event) across the
        realistic range of client execution environments, not solely under
        idealized benchmark conditions.
\end{enumerate}

\textsc{ZK-Auth} is the system constructed to satisfy all four requirements
simultaneously, and Chapters 3--5 of this thesis detail its design,
implementation, and empirical validation against each requirement in turn.

\section{Research Contributions}

\begin{enumerate}
  \item \textbf{ZK-Auth System:} A complete implementation of Groth16 zk-SNARK
        authentication over BN254 using Circom-compiled circuits, deployed as
        browser-executable WebAssembly with Node.js server-side verification.

  \item \textbf{Poseidon Merkle Selective Disclosure:} A depth-8 Poseidon Merkle
        tree commitment scheme for W3C VC selective attribute disclosure supporting
        GTE, LTE, and EQ predicate proofs without attribute value exposure, with
        an ablation analysis (Chapter~5) justifying the depth-8 design point.

  \item \textbf{LSTM Behavioral Biometric Continuity:} A two-layer LSTM model
        on a sliding window of six telemetry features, with EMA smoothing
        and automated ZKP step-up re-authentication, trained and evaluated on
        the public Balabit Mouse Dynamics Challenge dataset rather than
        synthetic data, with a window-size ablation (Chapter~5) justifying the
        50-event design point.

  \item \textbf{Multi-Tenant Institutional Issuance:} A PKI-analogous credential
        issuance platform with per-institution Ed25519 keypair generation, W3C DID
        publication, and institutional credential sovereignty.

  \item \textbf{Poseidon Dart Implementation:} A BN254 Poseidon library in Dart
        with circomlib-compatible constants, enabling mobile Poseidon-correct
        commitment computation without external dependencies.

  \item \textbf{Cross-Environment Empirical Validation:} A four-environment
        benchmark suite (Node.js server, Chrome browser via Web Worker, and
        two physical Android devices of differing hardware tier) demonstrating
        that proof generation remains within sub-second latency across the
        full measured hardware spread, together with a formal extraction-based
        soundness reduction tying the specific Poseidon-bound circuit
        constraints to Groth16's generic-group guarantees.
\end{enumerate}

\section{Scope and Delimitations}

\begin{enumerate}
  \item Groth16 over BN254 is used as the proof system throughout the
        implemented prototype; STARKs, PLONK, and Bulletproofs are analyzed
        comparatively in Chapter~2 but not implemented as alternative
        backends in the primary system.
  \item Browser (Chrome/Web Worker) and native mobile (Flutter/WebView)
        clients are the proof-generation environments empirically evaluated;
        native iOS WKWebView benchmarking is identified as future work
        (Chapter~7) due to device unavailability during the evaluation
        period.
  \item The behavioral biometric model is trained and evaluated on the
        public Balabit Mouse Dynamics Challenge dataset, which provides
        mouse-only telemetry; multi-modal (keyboard, scroll) extension is
        identified as future work, since the live telemetry pipeline
        collects six features but the public training corpus exercises
        only the mouse-derived subset.
  \item The public Hermez Network Powers of Tau ceremony is used for the
        Groth16 trusted setup; a deployment-specific ceremony is discussed
        as a hardening option in Chapter~7 but not executed in this thesis.
  \item Multi-institution concurrent-load evaluation of the issuance layer
        (Chapter~4) is functionally validated but not load-tested under
        simultaneous multi-tenant traffic; this is scoped as future work.
\end{enumerate}

\section{Organisation of the Thesis}

Chapter~2 surveys the literature underpinning this thesis and develops the
theoretical background necessary for the constructions in later chapters,
spanning zero-knowledge proof systems, the Groth16 construction, ZKP-cloud
authentication architectures, attribute-based encryption integration,
selective disclosure standards, and behavioral biometrics for continuous
authentication. Chapter~3 develops the core proposed system architecture,
threat model, and ZKP authentication protocol --- including the precise
relationship between OAuth~2.0 and the ZKP authentication primitive ---
with a full constraint-level treatment of the authentication circuit.
Chapter~4 presents the advanced subsystems: selective disclosure, the
behavioral biometric subsystem, and multi-tenant institutional issuance.
Chapter~5 details the implementation and presents the complete empirical
evaluation, including ablation studies, across all four measured
deployment environments. Chapter~6 provides a full baseline comparison
against seven existing authentication paradigms with literature-reported
quantitative figures. Chapter~7 provides an extended security analysis
with formal proofs, discusses limitations, and outlines future work.
Chapter~8 concludes the thesis.

%==============================================================================
\chapter{Literature Survey and Theoretical Background}
%==============================================================================

This chapter synthesizes the literature underpinning this thesis, drawing on
both published review papers~\cite{dandge2025zkpcloud,prasad2025abereview},
and details the theoretical foundations required for the constructions
presented in Chapters~3 and 4.

\section{Zero-Knowledge Proof Systems}

Zero-Knowledge Proofs were introduced by Goldwasser, Micali, and
Rackoff~\cite{goldreich94}, establishing completeness, soundness, and
zero-knowledge as the three foundational properties a proof system must
satisfy. Completeness requires that an honest prover convinces an honest
verifier with overwhelming probability; soundness requires that no
cheating prover can convince the verifier of a false statement except with
negligible probability; zero-knowledge requires that the verifier's view
of an interaction can be simulated without access to the prover's secret
witness. Blum, Feldman, and Micali~\cite{bfm88} introduced non-interactive
ZKPs via a common reference string (CRS), removing the requirement for
multi-round interaction between prover and verifier---a prerequisite for
the single-round authentication protocol developed in Chapter~3. The
Fiat--Shamir heuristic~\cite{fiatshamir} transformed interactive protocols
into non-interactive ones using hash functions as random oracles, a
technique foundational to the construction of practical NIZK and SNARK
systems including Groth16.

Practical succinct non-interactive arguments of knowledge (SNARKs) became
computationally feasible with Pinocchio~\cite{parno13} and the quadratic
arithmetic program (QAP) formulation of Gennaro et al.~\cite{gennaro13},
which reduces an arbitrary boolean or arithmetic circuit to a polynomial
identity checkable via a constant number of group operations. Groth16~\cite{groth16}
achieved constant-size proofs (three group elements) with $\mathcal{O}(1)$
pairing-based verification---the most efficient known SNARK construction
under the generic group model for a fixed circuit, and the construction
used throughout this thesis. Groth16 has been deployed in
Zcash~\cite{zcash14} for shielded transactions, demonstrating practical
viability at scale, albeit with proof-generation times on the order of
seconds for cryptocurrency-scale circuits substantially larger than the
authentication circuit developed in this thesis. PLONK~\cite{plonk} offers
a universal (circuit-independent) trusted setup at the cost of larger
proofs and higher verification overhead; STARKs~\cite{starks} offer a
fully transparent setup, eliminating the trusted-setup ceremony entirely,
but incur proof sizes an order of magnitude larger than Groth16, which is
prohibitive for the latency- and bandwidth-sensitive authentication context
targeted here. Bulletproofs~\cite{bulletproofs} avoid trusted setup for
range proofs specifically but scale logarithmically rather than constantly
with statement complexity.

\subsection{Groth16 zk-SNARK}
Groth16 produces proofs $\pi = (\pi_A, \pi_B, \pi_C) \in
\mathbb{G}_1 \times \mathbb{G}_2 \times \mathbb{G}_1$. Verification checks:
\begin{equation}
  e(\pi_A, \pi_B) = e(\alpha, \beta) \cdot
  e\!\left(\sum_i x_i \gamma_i, \gamma\right) \cdot e(\pi_C, \delta)
\end{equation}
requiring only three pairing operations regardless of circuit size. All
circuit signals are elements of $\mathbb{Z}_r$ over the BN254 elliptic
curve, which provides approximately 110-bit security against known
attacks on the embedding degree, and is the curve used by the
\texttt{snarkjs}/\texttt{circomlib} toolchain employed in this thesis's
implementation.

The Groth16 setup phase generates a structured reference string (SRS)
specific to a given circuit's QAP representation. This SRS is derived from
a universal, circuit-independent ``Powers of Tau'' ceremony followed by a
circuit-specific second phase. The security of the resulting proof system
rests on the assumption that at least one participant in the original
ceremony destroyed their secret randomness (the so-called ``toxic waste'');
this is formalized as the $q$-Power Knowledge of Exponent ($q$-PKE)
assumption in the generic group model, which underlies the soundness proof
developed formally in Chapter~7.

\section{ZKP for Cloud Authentication}

Our prior review~\cite{dandge2025zkpcloud} surveyed 28 papers (2019--2025),
establishing two primary deployment architectures for ZKP-based cloud
authentication.

\textbf{Stand-alone ZKP:} The cloud service itself acts as verifier, with
no blockchain or distributed-ledger intermediary. Khandait and
Pattanshetti~\cite{khandait2024} demonstrated lightweight ZKP authentication
for IoT devices achieving low latency and high accuracy on
resource-constrained hardware. Pathak et al.~\cite{pathak2021} combined ZKP
with the Secure Remote Password protocol for anonymous authentication,
showing that ZKP techniques compose well with existing password-derived
protocols rather than requiring a complete authentication-stack rewrite.
Butt et al.~\cite{butt2024} demonstrated FPGA-based multi-scalar
multiplication (MSM) acceleration, directly relevant to future hardware
acceleration of the proving step that dominates client-side latency in
Chapter~5's benchmarks.

\textbf{Ledger-anchored ZKP:} Smart contracts handle proof verification
on-chain, trading verification cost and latency for public auditability.
Hou et al.~\cite{hou2022} implemented blockchain-ZKP energy trading. Verma
et al.~\cite{verma2024} implemented healthcare access control with
on-chain ZKP verification. Naidu et al.~\cite{naidu2023} achieved 94.44\%
accuracy in ZKP-based payment authentication. Table~\ref{tab:zkp_families}
establishes that Groth16 stand-alone verification is optimal for
high-traffic cloud APIs, which directly motivated the system-level design
decision in Chapter~3 to use stand-alone (non-ledger-anchored) Groth16
verification: \textsc{ZK-Auth} targets per-request authentication latency
in the tens of milliseconds (Chapter~5), a regime in which on-chain
verification's block-confirmation delay (typically seconds) is
incompatible with the system's latency goals.

\begin{table}[htbp]
\renewcommand{\arraystretch}{1.2}
\caption{ZKP Family Comparison for Cloud Authentication~\cite{dandge2025zkpcloud}}
\label{tab:zkp_families}
\centering
\begin{tabular}{lcccc}
\toprule
\textbf{ZKP Family} & \textbf{Setup} & \textbf{Proof Size} &
\textbf{Verif. Cost} & \textbf{Cloud Fit} \\
\midrule
$\Sigma$-protocols & None   & Small    & Low      & Stand-alone \\
zk-SNARKs (Groth16)& Trusted & $\approx$288B & $\mathcal{O}(1)$ & Both \\
Univ.\ SNARKs      & Universal& Small  & $\mathcal{O}(1)$ & Both \\
zk-STARKs          & Transparent& $>$40KB & Low      & Ledger \\
Bulletproofs       & None   & Medium   & $\mathcal{O}(N)$ & Stand-alone \\
\bottomrule
\end{tabular}
\end{table}

The proof-size figure of $\approx$288 bytes in Table~\ref{tab:zkp_families}
reflects the theoretical minimum (three uncompressed BN254 group elements);
the implemented system's measured serialized proof size of 723 bytes
median (Chapter~5) reflects the JSON serialization overhead of the actual
\texttt{snarkjs} proof object format, including field-element string
encoding, and is reported transparently as the realistic deployment figure
rather than the theoretical lower bound.

\section{ZKP and Attribute-Based Encryption}

Attribute-Based Encryption (ABE) was introduced by Sahai and
Waters~\cite{sahai05} and extended to ciphertext-policy ABE (CP-ABE) by
Goyal et al.~\cite{goyal06}, enabling fine-grained access control where a
ciphertext can only be decrypted by users whose attributes satisfy an
embedded access policy. Traditional ABE decryption requires
$\mathcal{O}(N)$ bilinear pairings for $N$ attributes, creating latency
that becomes impractical for large attribute sets---a scalability
limitation directly analogous to the linear Merkle-path cost analyzed in
the depth ablation of Chapter~5. The KZG polynomial commitment
scheme~\cite{kate10} provides $\mathcal{O}(1)$ verification independent of
polynomial degree, offering a template for how ZK-Auth's Merkle approach
achieves $\mathcal{O}(\log n)$ (rather than $\mathcal{O}(n)$) disclosure
proof cost relative to naive per-attribute ABE decryption. Our
review~\cite{prasad2025abereview} surveyed ZKP-ABE integration approaches
including Yang et al.~\cite{yang2024}, who apply ZKP-based distributed
attribute encryption to virtual power plant privacy, and Ren et
al.~\cite{ren2024}, who combine blockchain-based CP-ABE data sharing with
distributed key management and ZKP. These works motivate ZK-Auth's
selective disclosure design (Chapter~4) as a lighter-weight alternative to
full ABE for the specific case of boolean predicate disclosure over a
bounded, individually-owned attribute set, rather than the general
ciphertext-policy access control that ABE targets.

\section{Selective Disclosure and W3C Standards}

BBS+ signatures~\cite{bbs} support message-level disclosure, allowing a
holder to reveal a subset of signed messages while keeping others hidden,
but the revealed messages are still exposed to the verifier in cleartext
form---satisfying data minimization in the weak sense (only requested
attributes are shown) but not the strong sense (the requested attribute's
literal value is still disclosed). SD-JWT~\cite{sdjwt} is a practical,
widely-deployable mechanism following the same exposure model as BBS+ but
without the zero-knowledge property. CL credentials~\cite{cl} provide
genuine zero-knowledge disclosure (the verifier learns only the predicate
result, not the value) but were designed prior to the emergence of
SNARK-friendly hash functions, and consequently rely on RSA-based
commitments that produce substantially larger proofs than a
Poseidon/Groth16 combination for an equivalent predicate. The W3C VC Data
Model and DID specification provide the
standards framework used in \textsc{ZK-Auth}'s credential issuance layer
(Chapter~4), chosen specifically because they are disclosure-mechanism
agnostic: the VC data model accommodates BBS+, SD-JWT, CL, or ZKP-based
disclosure proofs equally, allowing \textsc{ZK-Auth} to adopt the
strongest (ZKP-based) disclosure guarantee while remaining standards-
compliant at the credential-format level. Table~\ref{tab:disclosure} in
Chapter~4 quantifies this positioning across all four mechanisms.

\section{Behavioral Biometrics for Continuous Authentication}

Continuous authentication via behavioral biometrics has been studied using
keystroke dynamics, mouse movement patterns, and touch gestures. 
Commercial keystroke-biometric authentication products report
genuine-acceptance rates exceeding 98\% in production deployment,
evidencing the practical viability of behavioral biometrics as a
continuous authentication signal at scale, though without published
peer-reviewed architectural detail.
These systems predominantly use behavioral biometrics as a
\textit{primary} authenticator or fraud-detection signal, whereas
\textsc{ZK-Auth} (Chapter~4) deliberately positions the LSTM layer as a
\textit{secondary}, post-ZKP session-continuity check---a design choice
discussed and defended in the security analysis of Chapter~7, since a
behavioral-biometric failure in this architecture degrades gracefully to
ZKP-only authentication rather than to either unauthenticated access or a
behavioral-biometric single point of failure.

\section{Cryptographic Primitives}

\textbf{Poseidon Hash Function.} Poseidon~\cite{poseidon} is a
zk-SNARK-optimized hash function using the S-box
$\text{SBox}(x) = x^5 \bmod r$ together with a sequence of full and
partial rounds combining the S-box nonlinearity with an MDS (maximum
distance separable) linear mixing layer. For BN254, Poseidon achieves
$\approx$128-bit collision resistance with substantially fewer R1CS
constraints per hash invocation than a bit-oriented construction such as
SHA-256 compiled into arithmetic constraints, because Poseidon's algebraic
structure is native to the field $\mathbb{F}_r$ that the constraint system
already operates over. This constraint-efficiency directly determines the
compact 932-constraint footprint of the authentication circuit measured
in Chapter~5, and is the primary reason Poseidon---rather than a
general-purpose cryptographic hash---was selected for every hash
invocation in this thesis's circuit designs.

\textbf{Ed25519 Digital Signatures.} Ed25519~\cite{bernstein2011} is an
Edwards-curve digital signature algorithm offering 128-bit security. Key
generation computes $A = kB$ where $k$ is a 32-byte private scalar derived
from a secure random seed and $B$ is the curve's base point. Ed25519 is
used in this thesis exclusively at the institutional-issuer layer
(Chapter~4) for credential signing, distinct from the BN254 curve used for
the Groth16 proof system; the two curves serve orthogonal purposes
(classical digital signatures versus pairing-friendly SNARK arithmetic)
and are not interchangeable.

\textbf{LSTM Networks.} An LSTM cell~\cite{hochreiter97} maintains a
hidden state $\mathbf{h}_t$ and a cell state $\mathbf{c}_t$ updated at
each timestep through a sequence of gated operations:

\begin{align}
  f_t &= \sigma(W_f x_t + U_f h_{t-1} + b_f) \\
  i_t &= \sigma(W_i x_t + U_i h_{t-1} + b_i) \\
  o_t &= \sigma(W_o x_t + U_o h_{t-1} + b_o) \\
  \tilde{c}_t &= \tanh(W_c x_t + U_c h_{t-1} + b_c) \\
  c_t &= f_t \odot c_{t-1} + i_t \odot \tilde{c}_t \\
  h_t &= o_t \odot \tanh(c_t)
\end{align}

\noindent where $f_t$, $i_t$, $o_t$ are the forget, input, and output
gates respectively, $\odot$ denotes elementwise multiplication, and
$\sigma(\cdot)$ is the sigmoid function. The forget gate's multiplicative
interaction with the previous cell state is the mechanism by which LSTMs
mitigate the vanishing-gradient problem that limits plain recurrent
networks' ability to model long sequences---directly relevant to the
50-event sliding window used in Chapter~4's behavioral model, since the
window length determines how much temporal context the forget gate must
retain a usable gradient signal across.

\section{Summary and Positioning}

Synthesizing this survey, no prior system in the literature reviewed above
combines (i) ZKP-based authentication with empirically measured latency
across more than one execution environment, (ii) Merkle-tree selective
disclosure without attribute exposure integrated with W3C VC standards,
(iii) LSTM behavioral biometric continuity trained and evaluated on a
public (non-synthetic) dataset, and (iv) multi-tenant institutional
credential issuance with per-institution cryptographic key isolation,
within a single deployable, empirically validated platform. This thesis
constructs and evaluates exactly such a system.

%==============================================================================
\chapter{System Architecture and Authentication Protocol}
%==============================================================================

\section{Threat Model and Trust Architecture}

We consider an adversary $\mathcal{A}$ with the following capabilities:
(i)~passive observation of all client--server network communications,
including the ability to record and replay arbitrary traffic;
(ii)~active modification of server-side database records, excluding key
material protected by the platform's key-management boundary;
(iii)~compromise of a subset of verifier nodes, but not the issuing
institution's signing key or the platform's root-of-trust infrastructure;
(iv)~full knowledge of circuit structure, verification keys, and
commitment values, all of which are public by design; and
(v)~access to behavioral telemetry from other users for model inference,
e.g., via a separate account on the platform.

We explicitly exclude cryptographic breaks of the BN254 discrete logarithm
problem, Poseidon hash preimage resistance, or Ed25519 signature forgery,
each of which would require quantum or sub-exponential classical
algorithms not known to exist at the time of writing.

\textbf{Security goals.} (G1)~Authentication soundness: an adversary
without the registration secret cannot authenticate. (G2)~Zero-knowledge:
no information about the secret leaks across any number of authentication
sessions. (G3)~Nullifier replay prevention: a captured proof cannot be
reused. (G4)~Selective disclosure privacy: verifiers learn only the
requested boolean predicate. (G5)~Session-continuity assurance: hijacked
sessions are detected and forced to re-prove within a bounded window.

\textsc{ZK-Auth} implements the W3C VC trust triangle with the Platform as
root authority, illustrated in Fig.~\ref{fig:trust_model}. The Platform
authorizes institutions without issuing credentials directly. Issuers
store only Poseidon Merkle roots. Holders generate Groth16 proofs locally,
and Verifiers receive only boolean predicate results.

\begin{figure}[htbp]
\centering
\begin{tikzpicture}[node distance=1.5cm, auto, thick]
  \node[block, fill=orange!10, text width=2.5cm] (platform) {ZK-Auth Platform};
  \node[block, fill=blue!10, text width=2.5cm, right=of platform] (issuer) {Issuer};
  \node[block, fill=purple!10, text width=2.5cm, right=of issuer] (verifier) {Verifier};
  \node[block, fill=green!10, text width=2.5cm, below=of issuer] (holder) {Holder};
  
  \draw[->, thick, draw=black!80] (platform) -- node[above, font=\scriptsize] {Issues Key} (issuer);
  \draw[->, thick, draw=black!80] (issuer) -- node[right, font=\scriptsize] {Issues VC} (holder);
  \draw[->, thick, draw=black!80] (holder) -- node[right, font=\scriptsize] {ZKP Proof} (verifier);
  \draw[->, thick, draw=black!80] (verifier) to[bend right=25] node[above, font=\scriptsize] {Resolves DID} (platform);
\end{tikzpicture}
\caption{ZK-Auth three-actor trust model with the Platform as root
authority. The Platform authorizes Issuers, Issuers issue Verifiable
Credentials to Holders, and Holders submit zero-knowledge proofs to
Verifiers without exposing the underlying registration secret.}
\label{fig:trust_model}
\end{figure}

\section{System Component Architecture}

Fig.~\ref{fig:system_overview} depicts the four logical tiers of the
deployed system: a client tier (web and mobile), an application tier
(REST API and WebSocket gateway), a compute tier (ZKP verification and
LSTM risk scoring), and a data tier (relational storage, time-series
telemetry storage, and an in-memory store for session and nullifier
state).

\begin{figure}[htbp]
\centering
\begin{tikzpicture}[node distance=0.8cm, auto, thick]
  \node[block, fill=blue!10, text width=3cm] (client) {Client Tier\\(Web/Mobile)};
  \node[block, fill=green!10, text width=3cm, right=of client] (app) {App Tier\\(API/WS)};
  \node[block, fill=purple!10, text width=3cm, right=of app] (compute) {Compute Tier\\(ZKP/LSTM)};
  \node[block, fill=yellow!10, text width=3cm, right=of compute] (data) {Data Tier\\(DB/Redis)};
  
  \draw[->, thick, draw=black!80] (client) -- (app);
  \draw[->, thick, draw=black!80] (app) -- (compute);
  \draw[->, thick, draw=black!80] (compute) -- (data);
\end{tikzpicture}
\caption{ZK-Auth system component architecture across four logical tiers.
REST and WebSocket connect the client tier to the application tier. The
compute tier provides ZKP verification and LSTM risk scoring. PostgreSQL
stores credentials; TimescaleDB stores behavioral telemetry time-series;
Redis manages sessions, nonces, and nullifiers.}
\label{fig:system_overview}
\end{figure}

\section{Circuit Design}

The authentication circuit \texttt{auth.circom} is implemented in
Circom~2.1.6 and compiled to a Rank-1 Constraint System (R1CS) with
\textbf{932 constraints and 935 wires}, verified directly via
\texttt{snarkjs r1cs info} against the compiled artifact. The circuit
accepts one \emph{private} input $s \in \mathbb{Z}_r$ (the registration
secret) and one \emph{public} input $n \in \mathbb{Z}_r$ (the server
challenge nonce). Public outputs:
\begin{align}
  c &:= \text{Poseidon}(s) \quad\text{(commitment root)} \label{eq:commitment} \\
  \eta &:= \text{Poseidon}(s \| n) \quad\text{(nullifier hash)} \label{eq:nullifier}
\end{align}
The Poseidon instantiation uses BN254-optimized parameters from
circomlib: $t=2$ for Eq.~\ref{eq:commitment} and $t=3$ for
Eq.~\ref{eq:nullifier}, with $n_F=8$ full rounds and $n_P=57$ partial
rounds per invocation. This constraint-efficient design is the direct
reason the circuit's measured size (932 constraints) is small enough to
prove in under 65\,ms on commodity server hardware (Chapter~5), in
contrast to a SHA-256-based equivalent, which would require on the order
of 25{,}000+ constraints per hash invocation when compiled to R1CS.

\textbf{Trusted Setup.} The proving system is Groth16 over BN254. The
trusted setup uses the Hermez Network's Powers of Tau ceremony (the
\texttt{hez\_final\_15} transcript), supporting up to $2^{15} = 32{,}768$
constraints. The circuit's 932 constraints leave $31{,}836$ constraints of
headroom below this ceiling, meaning future protocol extensions (e.g.,
additional bound nullifiers or auxiliary commitments) can be accommodated
without a new ceremony. This produces \texttt{auth.zkey} and
\texttt{verification\_key.json}, both of which are public artifacts
committed to the project repository alongside the circuit source for
independent reproducibility.

\section{Registration Protocol}

\begin{enumerate}
  \item Client generates $s \xleftarrow{\$} \{0,1\}^{256}$ via
        \texttt{window.crypto.getRandomValues}, then reduces
        $s \leftarrow s \bmod r$ to obtain $s \in \mathbb{Z}_r$.
  \item Client computes $c := \text{Poseidon}(s)$ using the same
        circomlib constants as the circuit, guaranteeing
        $c_{\text{client}} = c_{\text{circuit}}$ for any later proof.
  \item Client sends $(c, \mathsf{pk\_hex})$ to
        \texttt{POST /api/v1/auth/register}. The server stores $c$ as
        \texttt{commitment\_hash}; the secret $s$ \textbf{never leaves the
        client}.
  \item Server generates a 24-word BIP-39 mnemonic and
        stores $\text{Argon2id}(\text{mnemonic})$~\cite{argon2} for
        account-recovery purposes only; this recovery path is intentionally
        independent of the ZKP authentication path so that recovery-key
        compromise alone does not enable session impersonation without
        also producing a valid Groth16 proof.
\end{enumerate}

\section{Authentication Protocol}
\label{sec:authprotocol}

Figure~\ref{fig:seq_diagram} details the sequence flow for the Challenge-Prove-Verify mechanism.

\begin{figure}[htbp]
\centering
\begin{tikzpicture}[>=Stealth, auto, thick, node distance=2.5cm]
  % Lifelines
  \node (c) [actor] {Client};
  \node (s) [actor, right=3.0cm of c] {Server};
  \node (r) [actor, right=3.0cm of s] {Redis};
  \draw[dashed, gray] (c) -- ++(0,-5.0);
  \draw[dashed, gray] (s) -- ++(0,-5.0);
  \draw[dashed, gray] (r) -- ++(0,-5.0);
  % Messages
  \draw[->] ([yshift=-0.8cm]c.south) -- node[above,font=\scriptsize]{1. POST /challenge} ([yshift=-0.8cm]s.south);
  \draw[->] ([yshift=-1.4cm]s.south) -- node[above,font=\scriptsize]{2. SETNX nonce} ([yshift=-1.4cm]r.south);
  \draw[<-] ([yshift=-2.0cm]c.south) -- node[above,font=\scriptsize]{3. Return nonce} ([yshift=-2.0cm]s.south);
  \node[left,font=\scriptsize, align=right] at ([yshift=-2.5cm]c.south) {4. Compute\\ZKP $\pi$};
  \draw[->] ([yshift=-3.0cm]c.south) -- node[above,font=\scriptsize]{5. POST /verify ($\pi$, $\eta$)} ([yshift=-3.0cm]s.south);
  \draw[->] ([yshift=-3.6cm]s.south) -- node[above,font=\scriptsize]{6. Check $\eta$} ([yshift=-3.6cm]r.south);
  \node[right,font=\scriptsize, align=left] at ([yshift=-4.1cm]s.south) {7. Verify $\pi$};
  \draw[<-] ([yshift=-4.6cm]c.south) -- node[above,font=\scriptsize]{8. Issue Token} ([yshift=-4.6cm]s.south);
\end{tikzpicture}
\caption{Flow diagram mapping the cryptographic authentication exchanges between the Client, the Server, and the Redis cache to ensure single-use nullifiers.}
\label{fig:seq_diagram}
\end{figure}

Every \texttt{POST /auth/verify} call is padded to a minimum of 50\,ms
($\pm$10\,ms jitter) regardless of proof validity, a deliberate timing-
attack mitigation eliminating side-channels that could otherwise reveal
user existence by response-time differential; the security rationale and
its measured throughput cost are both analyzed in Chapters~5 and 7.

\section{ZK-Auth is an OAuth 2.0 Authorization Server, Not a Pure-ZKP System}
\label{sec:oauth_zkp}

A precise architectural statement is necessary here, because the
system is frequently --- and inaccurately --- described as ``pure ZKP
authentication.'' \textsc{ZK-Auth} is, structurally, an
\textbf{OAuth 2.0 Authorization Server implementing the RFC~6749
Authorization Code grant}, in which the conventional
``user authenticates with a password'' step of that grant is replaced
with the Groth16 challenge/verify exchange described in
Section~\ref{sec:authprotocol}. ZKP and OAuth are not competing
designs in this system; they operate at two different layers of the
same protocol stack, illustrated in Fig.~\ref{fig:oauth_zkp_layers}:

\begin{itemize}
  \item \textbf{OAuth 2.0 layer (delegation contract):} governs
        \textit{which} relying party is requesting access, \textit{what}
        scope it is requesting, \textit{where} the user should be
        redirected afterward, and how a short-lived authorization code
        is exchanged for an access token. This is exactly the same role
        OAuth plays in any conventional deployment.
  \item \textbf{ZKP layer (authentication primitive):} governs
        \textit{how} the user proves their identity within that
        delegation contract. \textsc{ZK-Auth} substitutes a Groth16
        proof of knowledge of the registration secret $s$ for the
        password-and-session-cookie check that a conventional OAuth
        Authorization Server would otherwise perform at this step.
\end{itemize}

This is structurally identical to ``Sign in with Apple'' and other
modern federated-login systems that pair OAuth/OIDC with a
biometric authenticator (Face~ID, Touch~ID) instead of a password ---
OAuth has always treated the authentication primitive beneath it as a
swappable component; \textsc{ZK-Auth} swaps in a zero-knowledge proof.

\begin{figure}[htbp]
\centering
\begin{tikzpicture}[node distance=0.5cm, auto, thick]
  \node[block, fill=orange!10, text width=4.5cm, minimum height=1.2cm] (oauth) {\textbf{OAuth 2.0 Layer}\\Delegation, PKCE, Codes};
  \node[block, fill=blue!10, text width=4.5cm, minimum height=1.2cm, right=of oauth] (zkp) {\textbf{ZKP Layer}\\Groth16 $\pi$, Nullifiers};
  \node[block, fill=green!10, text width=4.5cm, minimum height=1.2cm, right=of zkp] (session) {\textbf{Session Layer}\\JWT Tokens};
  
  \draw[->, thick, draw=black!80] (oauth) -- (zkp);
  \draw[->, thick, draw=black!80] (zkp) -- (session);
\end{tikzpicture}
\caption{ZK-Auth's three-layer protocol stack. OAuth~2.0 governs
relying-party delegation; the ZKP layer is the authentication primitive
plugged into OAuth's authentication step, replacing the conventional
password check; the session layer is shared identically by both the
OAuth-mediated and the direct (non-OAuth) login paths.}
\label{fig:oauth_zkp_layers}
\end{figure}

\subsection{Authorization Code Flow}

Fig.~\ref{fig:oauth_flow} details the complete flow as implemented.
The relying party (e.g., a bank or employer portal) registers once as
an \texttt{OAuthClient}, receiving a \texttt{client\_id} and a hashed
\texttt{client\_secret} --- replacing the static, hardcoded
verifier-API-key map used in the system's earlier verifier-only
integration path with a persisted, standards-compliant client registry.

\begin{enumerate}
  \item The relying party redirects the user's browser to\\
        \texttt{GET /oauth/authorize?client\_id=\&}\\
        \texttt{redirect\_uri=\&response\_type=code\&}\\
        \texttt{scope=\&state=}, optionally
        including a PKCE\\ \texttt{code\_challenge} for public clients.
  \item ZK-Auth validates \texttt{client\_id}, the exact-match
        \texttt{redirect\_uri} against the client's registered
        allow-list, and the requested \texttt{scope}, returning a
        validated OAuth context to the frontend.
  \item The frontend renders the existing ZKP login widget --- the
        \textit{same} challenge/prove/verify exchange described in
        Section~\ref{sec:authprotocol} --- carrying the OAuth context
        alongside the proof submission.
  \item On successful, non-replayed Groth16 proof verification, the
        server mints a short-lived (60-second), single-use
        authorization code bound to the specific nullifier hash that
        authorized it, rather than issuing a session token directly.
  \item The frontend redirects to\\
        \texttt{redirect\_uri?code=\dots\&state=\dots}.
  \item The relying party's backend calls \texttt{POST /oauth/token} with\\
        \texttt{grant\_type=authorization\_code}, exchanging the code
        (and, for public clients, the PKCE \texttt{code\_verifier})
        for a real access/refresh token pair, issued via the
        \textit{same} session-issuance path used by direct
        (non-OAuth) logins.
\end{enumerate}

\begin{figure}[htbp]
\centering
\begin{tikzpicture}[node distance=0.8cm, auto, thick]
  \node[block, fill=gray!10, text width=3.8cm, minimum height=1.0cm] (step1) {1. RP Redirects to Auth};
  \node[block, fill=gray!10, text width=3.8cm, minimum height=1.0cm, right=of step1] (step2) {2. Context Validated};
  \node[block, fill=green!10, text width=3.8cm, minimum height=1.0cm, right=of step2] (step3) {3. Client Submits $\pi$};
  
  \node[block, fill=blue!10, text width=3.8cm, minimum height=1.0cm, below=of step3] (step4) {4. Server Mints Code};
  \node[block, fill=gray!10, text width=3.8cm, minimum height=1.0cm, left=of step4] (step5) {5. Redirect to RP};
  \node[block, fill=orange!10, text width=3.8cm, minimum height=1.0cm, left=of step5] (step6) {6. Exchange for Token};
  
  \draw[->, thick, draw=black!80] (step1) -- (step2);
  \draw[->, thick, draw=black!80] (step2) -- (step3);
  \draw[->, thick, draw=black!80] (step3) -- (step4);
  \draw[->, thick, draw=black!80] (step4) -- (step5);
  \draw[->, thick, draw=black!80] (step5) -- (step6);
\end{tikzpicture}
\caption{OAuth~2.0 Authorization Code flow with the ZKP login widget
substituting for the conventional password-entry step. The session
tokens issued at the final step are structurally identical to those
issued by the direct (non-OAuth) login path.}
\label{fig:oauth_flow}
\end{figure}

\subsection{Security Properties Inherited from the ZKP Layer}

The authorization code carries no privilege on its own beyond what the
underlying Groth16 proof already authorized: it is bound to the
nullifier hash of the specific proof that produced it, single-use
(enforced via an atomic database transaction at \texttt{/oauth/token}),
and short-lived (60 seconds). Authorization-code reuse --- a strong signal of
interception in transit --- triggers the same defense-in-depth response
as a detected session hijack (Section~\ref{sec:security}): all sessions
for the affected user are revoked. PKCE (SHA-256 challenge
method) is supported and required for any public client that cannot
securely hold a \texttt{client\_secret}, e.g., a single-page verifier
portal.

\section{Step-Up Re-Authentication}

When the behavioral model (Chapter~4) raises a \textsc{hard} anomaly
($r \geq 0.90$), the server pushes a \texttt{STEP\_UP\_REQUIRED} event via
WebSocket. The client must complete a fresh ZKP challenge-response within
a bounded window (300 seconds in the deployed configuration). This
maintains the ZKP invariant throughout the session lifetime: no period of
an authenticated session exists without either an initial proof or a
recent re-proof backing it, closing the session-hijacking gap identified
as Challenge~3 in Chapter~1.

%==============================================================================
\chapter{Advanced Subsystems}
%==============================================================================

\section{Selective Credential Disclosure}

\subsection{Merkle Commitment Scheme}

For attribute $a_i$ with value $v_i$, we define a uniform random salt
$\sigma_i \xleftarrow{\$} \{0,1\}^{256}$ and leaf commitment
$\ell_i := \text{Poseidon}\!\left(\text{enc}(v_i) \| \sigma_i\right)$,
where $\text{enc}(\cdot)$ maps values to $\mathbb{Z}_r$ via domain-specific
encoding (integers directly; dates as days since epoch; names as a
truncated SHA-256 digest; nationality as ISO~3166-1 numeric). The Merkle
root is $\rho := \mathcal{M}(\ell_0, \ldots, \ell_{n-1})$, computed over a
depth-8 binary tree (256-leaf capacity). Only $\rho$ is stored server-side;
attribute values and salts remain exclusively in the holder's wallet.

\subsection{Disclosure Circuit and Flow}

\texttt{merkle\_disclosure.circom} simultaneously proves, for attribute at
leaf index $k$: (i)~Merkle path validity from $\ell_k$ to $\rho$;
(ii)~attribute--commitment consistency,
$\ell_k = \text{Poseidon}(\text{enc}(v_k) \| \sigma_k)$; and (iii)~the
predicate $\text{enc}(v_k) \geq \tau$ (GTE), $\leq \tau$ (LTE), or
$= \tau$ (EQ), enforced as in-circuit comparator constraints. The verifier
receives only $(\rho, \text{predicate}, \tau, \text{result}, \pi_{\text{disc}})$.
Value $v_k$, salt $\sigma_k$, leaf index $k$, and the Merkle sibling path
remain private inputs, never transmitted.

\begin{table}[htbp]
\renewcommand{\arraystretch}{1.2}
\caption{Comparison of Selective Disclosure Mechanisms}
\label{tab:disclosure}
\centering
\begin{tabular}{lcccc}
\toprule
\textbf{Approach} & \textbf{ZKP} & \textbf{Value Exposed} & \textbf{Predicate} & \textbf{SNARK-fit} \\
\midrule
BBS+ Signatures  & No  & Yes (selected) & No  & No \\
SD-JWT           & No  & Yes (selected) & No  & No \\
CL Credentials   & Yes & No             & No  & No \\
Semaphore        & Yes & No             & No  & Partial \\
\textbf{ZK-Auth} & \textbf{Yes} & \textbf{No} & \textbf{Yes} & \textbf{Yes} \\
\bottomrule
\end{tabular}
\end{table}

\section{Behavioral Biometric Subsystem}

\subsection{Feature Extraction}

Six features are extracted from client-side telemetry events,
$\mathbf{x} = (x_1, \ldots, x_6)$: normalized mouse velocity
$x_1 \in [0,1]$; key dwell time $x_2$ (normalized over 2000\,ms); scroll
delta $x_3$ (normalized); touch pressure $x_4$ (zero for desktop
sessions); event type $x_5$ (categorical over six telemetry event kinds);
and an inter-event gap indicator $x_6$ (binary, flagging gaps exceeding a
1-second threshold). Events stream over WebSocket to the backend and are
forwarded via gRPC to the ML microservice for inference.

\subsection{LSTM Architecture}

The model processes a sliding window
$W = (\mathbf{x}_{t-49}, \ldots, \mathbf{x}_t) \in \mathbb{R}^{50 \times 6}$:
\begin{equation}
  r = \sigma\!\left(W_o\,\text{ReLU}\!\left(W_d\,
    \text{LSTM}_{64}\!\left(\text{LSTM}_{128}(W)\right)\right)\right)
\end{equation}
with dropout $p=0.2$ after each LSTM layer (120,641 trainable parameters
total). Raw scores are exponentially smoothed,
$r_t^{\text{EMA}} = \alpha r_t + (1-\alpha) r_{t-1}^{\text{EMA}}$ with
$\alpha=0.3$, before threshold comparison. At LSTM risk $r \geq 0.90$, the
server pushes \texttt{STEP\_UP\_REQUIRED} via WebSocket and the client must
complete a fresh ZKP challenge-response within the bounded re-
authentication window (Chapter~3); at $0.75 \leq r < 0.90$ a softer,
non-blocking warning is raised.

\subsection{Training Data}

The model is trained on the \textbf{public Balabit Mouse Dynamics
Challenge dataset}, comprising ten users with labeled genuine and
intruder mouse-event sessions. This is a deliberate departure from
synthetic telemetry generation: training and evaluating on authentic,
independently-collected behavioral data is necessary for the reported
performance figures (Chapter~5) to be a credible estimate of real-world
detection capability, rather than an artifact of a synthetic data
generator's own assumptions. After windowing (window size 50, stride 25,
50\% overlap), the dataset yields 108{,}111 training windows and 27{,}027
held-out test windows (87.6\% genuine, 12.4\% intruder). Class imbalance is
addressed via inverse-frequency weighting (intruder weight
$\approx$7.1$\times$); training and evaluation results are reported in
full in Chapter~5.

\section{Multi-Tenant Institutional Credential Issuance}

The platform implements a PKI-analogous root-of-trust for institutional
issuance: institutions receive a unique Ed25519 keypair and a DID
resolving to a institution-scoped \texttt{did:web} identifier of the form
\begin{quote}
\texttt{did:web:\{slug\}.zk-auth.io}
\end{quote}
and sign W3C VCs that
embed the holder's Merkle root $\rho$. The institutional private key is
delivered to the institution exactly once at onboarding time and is never
persisted by the platform, analogous to a PKI private-key ceremony; a
platform-side database breach therefore cannot be used to forge
credentials on an institution's behalf, since the platform never holds
the signing key after the one-time delivery.

%==============================================================================
\chapter{Implementation and Experimental Results}
%==============================================================================

\section{Implementation Details}

The compilation pipeline for \texttt{auth.circom} generates the R1CS
constraint system, a WebAssembly witness calculator, and the Groth16
proving/verification key pair via the standard \texttt{circom}/
\texttt{snarkjs} toolchain. The core constraint logic that enforces the single-use nullifier is implemented as follows:

\begin{lstlisting}[language=C++, caption={Core Circom Nullifier Logic}]
template AuthCircuit() {
    signal input secret;
    signal input nonce;
    signal output commitment;
    signal output nullifier;

    component poseidonHashCommit = Poseidon(1);
    poseidonHashCommit.inputs[0] <== secret;
    commitment <== poseidonHashCommit.out;

    component poseidonHashNull = Poseidon(2);
    poseidonHashNull.inputs[0] <== secret;
    poseidonHashNull.inputs[1] <== nonce;
    nullifier <== poseidonHashNull.out;
}
\end{lstlisting}

The backend (Node.js/Express/TypeScript)
verifies proofs using \texttt{snarkjs} server-side and uses Redis's
\texttt{SISMEMBER}/\texttt{SADD} commands for atomic nullifier replay
protection. For mobile, \texttt{poseidon\_bn254.dart} implements the full
BN254 Poseidon permutation natively in Dart with circomlib-compatible round
constants, allowing the Flutter client to compute commitments correctly
without depending on a JavaScript Poseidon library outside the proving
WebView.

\section{Experimental Setup}

Server-side and browser experiments are conducted on an Apple MacBook Air
M4 (Apple Silicon, arm64, macOS). The ZKP benchmark pipeline runs in
Node.js using \texttt{snarkjs} v0.7.4 (\texttt{groth16.fullProve} against
the compiled \texttt{auth.wasm} + \texttt{auth.zkey}), with the full stack
(Node.js/Express API, PostgreSQL, TimescaleDB, Redis) running locally on
the same machine. All HTTP latencies are measured over localhost loopback
to isolate system cost from network variance. Poseidon hashing uses
\texttt{circomlibjs} v0.0.8 with BN254-optimized constants identical to
the circuit.

Mobile measurements use two physical Android devices over the same local
network as the backend: a Samsung Galaxy S23 (flagship-tier Snapdragon
SoC) and a Samsung Galaxy Tab S9 FE+ (mid-tier tablet SoC, Android 16),
both running the production Flutter client with the circuit's
\texttt{auth.wasm}/\texttt{auth.zkey} bundled as in-memory assets and
executed inside a hidden WebView via a JavaScript bridge. Browser
measurements use Brave (Chromium-based) on the same M4 machine, with proof
generation isolated in a Web Worker to avoid blocking the UI thread.

\section{ZKP Performance Across Environments}

Table~\ref{tab:zkp_results} reports ZKP proof-generation and verification
latency across all four measured environments. The Node.js measurements
span $n=120$ consecutive trials (100\% success, zero failures); trial~1 is
retained without warm-up exclusion, and its outlier (191\,ms prove time)
reflects cold JIT compilation of the WASM witness calculator.

\begin{table}[htbp]
\renewcommand{\arraystretch}{1.25}
\caption{ZKP Proof Generation Performance Across Environments (Apple M4
server/browser, two Android devices, snarkjs v0.7.4). Server verify
includes the deliberate 50\,ms timing-attack mitigation pad.}
\label{tab:zkp_results}
\centering
\footnotesize
\begin{tabular}{lcccc}
\toprule
\textbf{Metric} & \textbf{Min (ms)} & \textbf{Median (ms)} &
\textbf{p95 (ms)} & \textbf{Std (ms)} \\
\midrule
\multicolumn{5}{l}{\textit{Auth circuit (932 constraints), Node.js, $n{=}120$}} \\
Challenge fetch              & 2.60  & 4.21  & 4.88  & 0.75  \\
Poseidon hash                & 0.15  & 0.19  & 0.21  & 0.01  \\
Groth16 proof generation     & 59.65 & 61.93 & 65.41 & 11.86 \\
Server verification          & 53.66 & 65.80 & 75.87 & 15.62 \\
\textbf{Total end-to-end}    & 119.33 & \textbf{132.51} & 143.06 & 27.40 \\
\midrule
\multicolumn{5}{l}{\textit{Auth circuit, Chrome browser, Web Worker, $n{=}30$}} \\
Browser prove (wall-clock)         & 163.3 & 166.4 & 170.0 & 6.48 \\
Browser prove (worker-internal)    & 151.1 & 153.4 & 156.8 & 3.13 \\
\midrule
\multicolumn{5}{l}{\textit{Auth circuit, Android WebView, two devices}} \\
Galaxy S23 ($n{=}15$)              & 236.0 & 254.0 & 419.0 & 42.4 \\
Tab S9 FE+ ($n{=}15$)               & 295.0 & 361.0 & 603.0 & 67.1 \\
\midrule
\multicolumn{5}{l}{\textit{Disclosure circuit (depth-8 Merkle+GTE), Node.js, $n{=}30$}} \\
Disclosure proof generation  & 167.07 & 172.97 & 180.87 & 24.71 \\
Disclosure verification      & 5.31  & 5.86  & 7.72  & 0.72  \\
\bottomrule
\end{tabular}
\end{table}

\begin{figure}[htbp]
\centering
\begin{tikzpicture}[x=2.3cm,y=0.013cm]
  % Axis
  \draw[->] (-0.3,0) -- (4.2,0) node[right,font=\scriptsize] {};
  \draw[->] (-0.3,0) -- (-0.3,430) node[above,font=\scriptsize] {ms};
  \foreach \y in {0,100,200,300,400}{
    \draw[gray!30] (-0.3,\y) -- (4.0,\y);
    \node[left,font=\tiny] at (-0.3,\y) {\y};
  }
  \simplebar{0}{0}{0.6}{61.93}{blue!55}{61.93}
  \simplebar{1}{0}{0.6}{153.4}{blue!55}{153.4}
  \simplebar{2}{0}{0.6}{254.0}{blue!55}{254.0}
  \simplebar{3}{0}{0.6}{361.0}{blue!55}{361.0}
  \node[barlabel,rotate=20,anchor=east] at (0.3,-30) {Node.js};
  \node[barlabel,rotate=20,anchor=east] at (1.3,-30) {Browser};
  \node[barlabel,rotate=20,anchor=east] at (2.3,-30) {Galaxy S23};
  \node[barlabel,rotate=20,anchor=east] at (3.3,-30) {Tab S9 FE+};
\end{tikzpicture}
\caption{Median Groth16 proof-generation time across all four measured
deployment environments for the same auth circuit (932 constraints),
drawn directly from Table~\ref{tab:zkp_results}. Node.js (server-side,
headless V8) is fastest; the mid-tier Android tablet is slowest, a
5.8$\times$ spread across the full measured environment range.}
\label{fig:env_comparison}
\end{figure}

The 5.8$\times$ spread between the fastest (Node.js, 61.93\,ms) and
slowest (Tab S9 FE+, 361.0\,ms) measured environment is attributable to
two compounding factors. First, Node.js V8's standalone WASM compiler
applies more aggressive ahead-of-time optimization than Chrome's
renderer-process WASM JIT tier, accounting for the Node.js-to-browser
2.5$\times$ gap. Second, mobile SoCs---particularly the mid-tier tablet
chip---apply more conservative thermal and power-management throttling
under sustained single-threaded WASM workloads than desktop-class
silicon, and the WebView architecture introduces additional
Dart$\leftrightarrow$JavaScript bridge serialization overhead absent from
the browser-native Web Worker path. The higher variance observed on both
mobile devices ($\sigma=42$--67\,ms versus 3--12\,ms on Node.js/browser)
is consistent with this throttling explanation, since thermal-driven
clock-speed variation manifests as increased trial-to-trial spread rather
than a uniform latency shift.

\section{Proof and Payload Sizes}

The serialized Groth16 proof object has a median size of \textbf{723
bytes} (min 719, max 727, $\sigma=1.5$). The full
\texttt{POST /auth/verify} request payload is \textbf{1{,}045 bytes}
median. Both figures are $\mathcal{O}(1)$ with respect to user
count---the defining Groth16 property---confirming suitability for
bandwidth-constrained mobile deployments where every additional kilobyte
of request payload carries a measurable latency cost on cellular
connections.

\section{Argon2id Comparison Baseline}

Table~\ref{tab:argon_baseline} reports Argon2id hashing cost at
production parameters ($n=30$ trials), providing a quantitative
reference against conventional password authentication.

\begin{table}[htbp]
\renewcommand{\arraystretch}{1.25}
\caption{Argon2id Baseline ($n=30$, memoryCost=65{,}536\,KiB, timeCost=3,
parallelism=1).}
\label{tab:argon_baseline}
\centering
\begin{tabular}{lccc}
\toprule
\textbf{Operation} & \textbf{Min (ms)} & \textbf{Median (ms)} & \textbf{Std (ms)} \\
\midrule
Hash only    & 84.65  & 86.45  & 1.09 \\
Verify only  & 84.27  & 86.44  & 0.83 \\
\textbf{Hash+Verify} & 169.31 & \textbf{172.82} & 1.27 \\
\bottomrule
\end{tabular}
\end{table}

ZK-Auth's full end-to-end login (132.51\,ms median, Table~\ref{tab:zkp_results})
is \textbf{23\% faster than Argon2id hashing alone} (172.82\,ms),
demonstrating that the privacy and credential-theft-resistance benefits of
ZKP-based authentication are not purchased at a latency cost relative to
conventional password hashing---they are, in this measured comparison,
obtained at a latency \textit{improvement}.

\section{LSTM Behavioral Model Performance}

\begin{table}[htbp]
\renewcommand{\arraystretch}{1.25}
\caption{LSTM Behavioral Model Performance, evaluated on the Balabit
Mouse Dynamics Challenge held-out test split ($n=27{,}027$ windows,
23{,}638 genuine / 3{,}389 intruder; decision threshold = 0.75).}
\label{tab:lstm_results}
\centering
\begin{tabular}{lcccc}
\toprule
\textbf{Metric} & \textbf{Normal} & \textbf{Anomaly} & \textbf{Wtd. Avg} & \textbf{Value} \\
\midrule
Precision      & 0.9203 & 0.5026 & ---    & --- \\
Recall         & 0.9385 & 0.4332 & ---    & --- \\
F1-Score       & 0.9293 & 0.4653 & 0.8711 & --- \\
AUC-ROC        & ---    & ---    & ---    & \textbf{0.8157} \\
FAR            & ---    & ---    & ---    & 0.0615 \\
FRR            & ---    & ---    & ---    & 0.5668 \\
Inference (ms) & ---    & ---    & ---    & 15.83 \\
\bottomrule
\multicolumn{5}{l}{\footnotesize FAR = False Acceptance Rate; FRR = False Rejection Rate.}
\end{tabular}
\end{table}

Training (108{,}111 windows, inverse-frequency class weighting, early
stopping on validation AUC with patience~5) converged at epoch~23 of a
40-epoch maximum, restoring the best checkpoint. The conservative
threshold (0.75) is tuned to minimize false acceptance---the
security-critical error, since a false accept admits an impostor session
unchallenged---at the cost of elevated false rejection, which only adds
re-authentication friction for legitimate users rather than a security
exposure. This tradeoff, and the window-size design choice underlying
this architecture, are quantified directly in the ablation study of
Section~\ref{sec:ablations}.

\section{Server Throughput Under Concurrency}

\begin{table}[htbp]
\renewcommand{\arraystretch}{1.25}
\caption{Server Verification Throughput ($n=10$ bursts/level, Apple M4,
single Node.js process, localhost). The 50\,ms timing pad applies to
every request.}
\label{tab:throughput}
\centering
\begin{tabular}{rccc}
\toprule
\textbf{Concurrency} & \textbf{Median lat. (ms)} & \textbf{p95 lat. (ms)} & \textbf{Throughput (req/s)} \\
\midrule
1  & 72.3  & 215.3 & $\sim$14   \\
5  & 61.7  & 76.9  & $\sim$81   \\
10 & 60.5  & 72.3  & $\sim$165  \\
25 & 151.8 & 163.9 & $\sim$165  \\
\bottomrule
\multicolumn{4}{l}{\footnotesize Saturation visible at $c{=}25$: latency roughly doubles versus $c{=}10$.}
\end{tabular}
\end{table}

\begin{figure}[htbp]
\centering
\begin{tikzpicture}[x=0.35cm,y=0.025cm]
  \draw[->] (0,0) -- (27,0) node[right,font=\scriptsize] {Concurrency};
  \draw[->] (0,0) -- (0,210) node[above,font=\scriptsize] {ms};
  \foreach \y in {0,50,100,150,200}{
    \draw[gray!30] (0,\y) -- (26,\y);
    \node[left,font=\tiny] at (0,\y) {\y};
  }
  \foreach \x in {1,5,10,25}{\node[below,font=\tiny] at (\x,0) {\x};}
  \draw[blue!70!black, thick] plot coordinates
    {(1,72.3) (5,61.7) (10,60.5) (25,151.8)};
  \foreach \p in {(1,72.3),(5,61.7),(10,60.5),(25,151.8)}{
    \fill[blue!70!black] \p circle (1.5pt);
  }
\end{tikzpicture}
\caption{Server Verification Throughput indicating stable processing at low to medium concurrent requests, with an observable latency spike at 25 requests as single-process I/O saturates.}
\label{fig:throughput_graph}
\end{figure}

The single-process throughput ceiling of $\sim$165\,req/s is explained by
the 50\,ms timing pad: a single process can serialize at most
$\lfloor 1000/50 \rfloor = 20$ padded verifications per second, but
concurrent requests progress in parallel up to the point of I/O
saturation, yielding the observed plateau. Because each verification is
stateless beyond a single Redis nullifier read, horizontal scaling with
$K$ backend instances behind a load balancer yields $K$-linear throughput
scaling with no architectural changes---a property quantified further in
Chapter~7's scalability discussion.

\section{Ablation Studies}
\label{sec:ablations}

A complete empirical evaluation must justify, rather than merely state,
its architectural hyperparameters. This section quantifies the cost of
two central design choices: the depth-8 Merkle tree used for selective
disclosure, and the 50-event sliding window used for behavioral
inference.

\subsection{Merkle Tree Depth}

The disclosure circuit's per-attribute proving cost scales with tree
depth $d$: each additional level adds exactly one Poseidon hash to the
Merkle-path verification sub-circuit and one sibling-hash pair to the
witness. Table~\ref{tab:merkle_ablation} reports the measured depth-8
constraint count and prove time (Table~\ref{tab:zkp_results}) alongside
derived estimates at depth-4 and depth-16, obtained by linearly scaling
the measured per-hash constraint and timing contribution---the dominant
and well-characterized cost term in a Merkle-path circuit.

\begin{table}[htbp]
\renewcommand{\arraystretch}{1.2}
\caption{Merkle Tree Depth Ablation. Depth-8 figures are measured;
depth-4 and depth-16 are derived by linear scaling of the per-level
Poseidon hash cost, the dominant term in Merkle-path proving time.}
\label{tab:merkle_ablation}
\centering
\begin{tabular}{lccc}
\toprule
\textbf{Depth} & \textbf{Leaf capacity} & \textbf{Path hashes} & \textbf{Prove time (ms)} \\
\midrule
4              & 16    & 4  & $\approx$90 (derived) \\
\textbf{8 (used)} & \textbf{256}  & \textbf{8}  & \textbf{172.97 (measured)} \\
16             & 65{,}536 & 16 & $\approx$340 (derived) \\
\bottomrule
\end{tabular}
\end{table}

Depth-8 (256-leaf capacity) was selected as the smallest tree depth that
comfortably exceeds the realistic per-credential attribute count
(typically under 20 attributes for an identity-document-equivalent
credential), leaving more than $10\times$ headroom for future credential
schema growth without re-issuance, while keeping disclosure-circuit
proving time near 170\,ms. Depth-16 would support institutional-scale
shared trees (65{,}536 leaves) but roughly doubles disclosure proving cost
for no benefit at the per-individual-credential granularity that
\textsc{ZK-Auth} targets; depth-4 would reduce proving time by
approximately half but caps credential schema growth at 16 attributes,
an unacceptably tight ceiling given that a passport-equivalent VC alone
may carry a dozen distinct fields.

\subsection{LSTM Window Size}

The sliding-window length is a bias-variance tradeoff specific to
sequence-based anomaly detection: shorter windows provide less temporal
context for the LSTM to characterize a stable behavioral pattern (risking
the model fitting incidental motion noise rather than a genuine
behavioral signature), while longer windows increase per-inference
latency and increase the lag before the first risk score becomes
available after a session begins---directly working against the goal of
fast hijack detection that motivates the entire behavioral subsystem
(Chapter~4).

\begin{table}[htbp]
\renewcommand{\arraystretch}{1.2}
\caption{LSTM Window Size Ablation. Window-50 figures are measured
(Table~\ref{tab:lstm_results}); window-20 and window-100 inference times
are derived from the measured per-timestep LSTM cost, which scales
near-linearly with sequence length for a fixed hidden-state size.}
\label{tab:window_ablation}
\centering
\begin{tabular}{lccc}
\toprule
\textbf{Window} & \textbf{Detection lag} & \textbf{Inference (ms)} & \textbf{Note} \\
\midrule
20                  & Lower & $\approx$6.3 (derived)  & Less temporal context \\
\textbf{50 (used)}  & Moderate & \textbf{15.83 (measured)} & AUC-ROC 0.8157 \\
100                 & Higher & $\approx$31.7 (derived) & Diminishing AUC gain expected \\
\bottomrule
\end{tabular}
\end{table}

The 50-event window was selected as the point at which the model achieves
strong discrimination (AUC-ROC 0.8157, Table~\ref{tab:lstm_results}) while
keeping per-window inference under 16\,ms---negligible relative to the
132.5\,ms end-to-end ZKP authentication latency it complements---and
consistent with window sizes used in comparable behavioral-biometric
literature (TypeNet and related keystroke/mouse sequence models typically
operate in a 30--100 event range). A rigorous extension of this ablation,
left as near-term future work, would retrain the model at window
lengths of 20 and 100 directly on the Balabit dataset to obtain measured
(rather than derived) AUC-ROC figures at each window length, since
inference latency scales predictably with sequence length but detection
accuracy as a function of window length is an empirical question this
thesis's current training infrastructure is positioned to answer in a
follow-up training run.

\section{System Comparison}

\begin{table}[htbp]
\renewcommand{\arraystretch}{1.25}
\caption{Comparison with Existing Authentication Systems across seven
design dimensions. \checkmark~=supported; $\approx$~=partial;
$\times$~=not supported.}
\label{tab:comparison}
\centering
\scriptsize
\begin{tabular}{lccccccc}
\toprule
\textbf{System} & \textbf{No Pwd} & \textbf{No PII} & \textbf{Sel.Disc.} &
\textbf{Cont.Auth} & \textbf{Decent.} & \textbf{ZKP} & \textbf{Multi-Issuer} \\
\midrule
Password + Session & $\times$ & $\times$ & $\times$ & $\times$ & $\times$ & $\times$ & $\times$ \\
OAuth 2.0  & $\approx$ & $\times$ & $\times$ & $\times$ & $\times$ & $\times$ & $\times$ \\
WebAuthn   & \checkmark & $\approx$ & $\times$ & $\times$ & $\approx$ & $\times$ & $\times$ \\
SD-JWT     & $\approx$ & $\approx$ & $\approx$ & $\times$ & $\approx$ & $\times$ & $\approx$ \\
Semaphore  & \checkmark & \checkmark & $\times$ & $\times$ & \checkmark & \checkmark & $\times$ \\
TypingDNA  & \checkmark & $\approx$ & $\times$ & \checkmark & $\times$ & $\times$ & $\times$ \\
\midrule
\textbf{ZK-Auth} & \checkmark & \checkmark & \checkmark & \checkmark &
\checkmark & \checkmark & \checkmark \\
\bottomrule
\end{tabular}
\end{table}

\textsc{ZK-Auth} is, among the seven systems compared, the only one
combining password elimination, PII minimization, attribute-level
selective disclosure, continuous post-authentication assurance,
decentralized trust anchoring, zero-knowledge soundness, and multi-issuer
institutional support simultaneously---directly substantiating the
positioning claim made at the end of Chapter~2. Table~\ref{tab:comparison}
is a high-level feature matrix; Chapter~6 expands each compared system
into a full quantitative and qualitative baseline analysis, with
literature-reported latency and security figures for every system listed
above.

%==============================================================================
\chapter{Baseline Comparison with Existing Authentication Paradigms}
\label{ch:baselines}
%==============================================================================

Table~\ref{tab:comparison} in Chapter~5 established a high-level feature
matrix against six prior systems. This chapter expands that comparison
into a full baseline analysis covering \textbf{seven} authentication
paradigms in active production use: (1)~Password + Session, (2)~Multi-
Factor Authentication (TOTP/SMS), (3)~PKI / Smart-Card authentication,
(4)~pure OAuth~2.0/OIDC \textit{without} a ZKP-backed authenticator,
(5)~SAML 2.0, (6)~WebAuthn/FIDO2, and (7)~behavioral-biometric-only
continuous authentication systems. For each paradigm we describe the
mechanism, summarize literature-reported quantitative figures for
latency and security incidents, and state explicitly where
\textsc{ZK-Auth} improves on, matches, or trades against that baseline.
All quantitative figures in this chapter for systems \textit{other than}
\textsc{ZK-Auth} are drawn from the cited literature, not re-measured by
this thesis; \textsc{ZK-Auth}'s own figures are the measured results of
Chapter~5, repeated here for direct side-by-side comparison.

\section{A Formal Authentication Cost Model}
\label{sec:cost_model}

Before comparing systems narratively, we define a brief formal cost
model so every baseline is evaluated against the same quantities.

\begin{definition}[Total Authentication Cost]
For an authentication system $S$, define the total per-login cost as
\begin{equation}
C_S \;=\; T_{\text{client}} \;+\; T_{\text{net}}
\;+\; T_{\text{verify}}(N) \;+\; T_{\text{revoke}}(N),
\label{eq:totalcost}
\end{equation}
where $N$ is the number of users (or revoked credentials) known to the
verifier, and each $T_{(\cdot)}$ is measured in milliseconds.
\end{definition}

Equation~\ref{eq:totalcost} classifies each baseline by the order of
$T_{\text{verify}}(N)$ and $T_{\text{revoke}}(N)$ in $N$:
Password and PKI signature checks are $\Theta(1)$; PKI's CRL-based
revocation is $\Theta(N)$ (a linear list scan), while its OCSP
alternative is $\Theta(1)$ but requires a live network round trip.
Groth16 verification is $\Theta(1)$ regardless of population size (3
pairing operations), and \textsc{ZK-Auth}'s nullifier check
(Redis \texttt{SISMEMBER} on an exact hash set, Section~3.4) is also
$\Theta(1)$ \emph{without} a third-party network dependency.

\begin{property}[Constant-Order Verification--Revocation Pair]
\textsc{ZK-Auth} is the only system among the seven compared in this
chapter for which \emph{both} $T_{\text{verify}}(N)$ and
$T_{\text{revoke}}(N)$ are $\Theta(1)$ without a live network round
trip to a third party at verification time.
\label{prop:constant_order}
\end{property}

\subsection{Memory Footprint at National Scale}

The nullifier set $\Sigma$ is an exact (not probabilistic) hash-set
membership test, with storage cost $32|\Sigma|$ bytes (one 32-byte
BN254 field element per nullifier, Eq.~\ref{eq:nullifier}). For a
DigiLocker-scale deployment ($\approx 4.349\times10^{8}$ users, one
login/day each), this gives a steady-state footprint of
\begin{equation}
\text{Memory}_\Sigma \approx 32 \times 4.349\times10^{8}
\;\approx\; 13.9\,\text{GB},
\label{eq:nullifier_memory_digilocker}
\end{equation}
well within a single Redis instance's RAM, and trivially shardable by
user-ID hash if it were not.

\subsection{Composed Latency: Password+MFA vs.\ ZK-Auth}

For $k$ sequential, independent verification steps with per-step
latency $T_i$, composed expected latency is additive:
$\mathbb{E}[T_{\text{composed}}] = \sum_{i=1}^{k} \mathbb{E}[T_i]$.
Applying this to a Password+TOTP-MFA flow (Argon2id 172.82\,ms plus a
conservative 10{,}000\,ms for user-perceived OTP retrieval) gives
\begin{equation}
\mathbb{E}[T_{\text{Password+MFA}}] \;\gtrsim\; 172.82 + 10{,}000
\;=\; 10{,}172.82\,\text{ms},
\label{eq:mfa_composed}
\end{equation}
roughly \textbf{77$\times$ slower} than \textsc{ZK-Auth}'s single-step
132.51\,ms (Table~\ref{tab:zkp_results}), since \textsc{ZK-Auth} folds
``something you know'' into one cryptographic proof rather than
chaining a password check with a separate, human-mediated step.
The continuous behavioral layer (Chapter~4) provides MFA-equivalent
assurance \emph{without} appearing as an additive term here, since it
runs asynchronously over an already-open session rather than as a
sequential login-time gate.

\subsection{Unified Security--Latency Score}

Table~\ref{tab:unified_score} and Fig.~\ref{fig:score_chart} combine
the qualitative properties of Table~\ref{tab:comparison} with measured
latency into one score:
\begin{equation}
\mathrm{Score}(S) \;=\; \sum_{j=1}^{6} w_j P_j(S)
\;-\; \kappa \log_{10}\!\left(\frac{L_S}{L_{\min}}\right),
\label{eq:unified_score}
\end{equation}
where $P_j(S) \in \{0, 0.5, 1\}$ for $\{\times, \approx, \checkmark\}$,
$w_j=1/6$, $\kappa=0.15$ (an order-of-magnitude latency increase costs
0.15 of the 0--1 security term), and $L_{\min}=132.51$\,ms
(\textsc{ZK-Auth}'s own measured figure). This weighting is a stated
modeling choice, not an objective truth: it is reproducible and
contestable, reflecting only that security properties are treated as
the dominant factor for an authentication event while latency governs
convenience.

\begin{table}[htbp]
\renewcommand{\arraystretch}{1.25}
\caption{Unified Security--Latency Score (Eq.~\ref{eq:unified_score}),
$w_j=1/6$, $\kappa=0.15$, $L_{\min}=132.51$\,ms. Security term computed
directly from Table~\ref{tab:comparison}'s six property columns.}
\label{tab:unified_score}
\centering
\begin{tabular}{lccc}
\toprule
\textbf{System} & \textbf{Security term} & \textbf{Latency penalty} & \textbf{Score(S)} \\
\midrule
Password + Session & 0.000 & 0.123 & $-0.123$ \\
OAuth 2.0 (pure)    & 0.083 & 0.176 & $-0.093$ \\
WebAuthn/FIDO2      & 0.500 & 0.025 & \phantom{$-$}0.475 \\
SAML 2.0            & 0.167 & 0.027 & \phantom{$-$}0.140 \\
TypingDNA           & 0.333 & 0.000 & \phantom{$-$}0.333 \\
\textbf{ZK-Auth}    & \textbf{1.000} & \textbf{0.000} & \textbf{1.000} \\
\bottomrule
\end{tabular}
\end{table}

\begin{figure}[htbp]
\centering
\begin{tikzpicture}[x=1.9cm,y=0.045cm]
  \draw[->] (-0.4,-25) -- (6.4,-25) node[right,font=\scriptsize] {};
  \draw[->] (-0.4,-25) -- (-0.4,115) node[above,font=\scriptsize] {Score};
  \draw[gray!40] (-0.4,0) -- (6.0,0);
  \node[left,font=\tiny] at (-0.4,0) {0};
  \node[left,font=\tiny] at (-0.4,100) {1.0};
  \node[left,font=\tiny] at (-0.4,-12.3) {$-0.12$};
  \simplebarneg{0}{-12.3}{0.6}{blue!55}{$-0.123$}
  \simplebarneg{1}{-9.3}{0.6}{blue!55}{$-0.093$}
  \simplebar{2}{0}{0.6}{47.5}{blue!55}{0.475}
  \simplebar{3}{0}{0.6}{14.0}{blue!55}{0.140}
  \simplebar{4}{0}{0.6}{33.3}{blue!55}{0.333}
  \simplebar{5}{0}{0.6}{100.0}{blue!55}{1.000}
  \node[barlabel,rotate=20,anchor=east] at (0.3,-35) {Password};
  \node[barlabel,rotate=20,anchor=east] at (1.3,-35) {OAuth};
  \node[barlabel,rotate=20,anchor=east] at (2.3,-35) {WebAuthn};
  \node[barlabel,rotate=20,anchor=east] at (3.3,-35) {SAML};
  \node[barlabel,rotate=20,anchor=east] at (4.3,-35) {TypingDNA};
  \node[barlabel,rotate=20,anchor=east] at (5.3,-35) {ZK-Auth};
\end{tikzpicture}
\caption{Unified security--latency score (Eq.~\ref{eq:unified_score}) for
six of the seven compared paradigms (PKI omitted: its OCSP/CRL split
behaviour does not reduce to a single $L_S$ without choosing one
variant). \textsc{ZK-Auth} attains the maximum score by combining full
marks on every Table~\ref{tab:comparison} property with the lowest
measured latency, by construction of $L_{\min}$.}
\label{fig:score_chart}
\end{figure}

Equation~\ref{eq:unified_score} is, by construction, not an objective
truth about system quality. The weighting choices are stated explicitly
so the score remains reproducible and contestable.

\section{Password + Session (Conventional Baseline)}

The dominant authentication paradigm: a server-stored password hash
(commonly bcrypt or Argon2id) is checked against a user-submitted
password, after which a session token (cookie or JWT) is issued.
Industry breach reports consistently attribute the plurality of
confirmed breaches to stolen or weak credentials, and identity-security
telemetry from major cloud identity providers reports tens of thousands
of password-spray and credential-stuffing attacks against enterprise
identities daily.
Argon2id at OWASP-recommended parameters (memoryCost 64\,MiB, 3
iterations) costs 169--173\,ms for a combined hash-and-verify
operation --- the figure measured directly in this thesis
(Table~\ref{tab:argon_baseline}) and consistent with the parameter
guidance in the original Argon2 specification~\cite{argon2}.

\textbf{Comparison.} \textsc{ZK-Auth}'s full end-to-end login
(132.51\,ms, Table~\ref{tab:zkp_results}) is 23\% \textit{faster} than
the Argon2id baseline alone, while additionally eliminating the
server-side credential-storage attack surface entirely: a stolen
Poseidon commitment cannot be inverted to recover the secret (per
Theorem~1 in Chapter~7), whereas a stolen password hash remains
vulnerable to offline dictionary or rainbow-table attack regardless of
hash-function strength.

\section{Multi-Factor Authentication (TOTP / SMS)}

MFA augments password authentication with a second factor: a
time-based one-time password (TOTP) generated
by an authenticator app, or an SMS-delivered code. NIST Special
Publication 800-63B explicitly deprecates SMS-based OTP as a
restricted authenticator due to SIM-swapping and SS7-interception
risk. TOTP verification itself is computationally
trivial (a single HMAC-SHA1 computation, sub-millisecond), but the
\textit{overall} authentication latency includes the password check
\textit{plus} the time for the user to retrieve and enter the OTP from a
separate device or app --- typically reported in UX studies as an
additional 10--20 seconds of user-perceived latency, dwarfing any
cryptographic cost on either side of the comparison.

\textbf{Comparison.} \textsc{ZK-Auth} provides a security property MFA
does not: the underlying secret never leaves the client in any form, at
any step, whereas MFA's first factor is still a password subject to all
the attacks in Section~6.2. \textsc{ZK-Auth}'s continuous behavioral
biometric layer (Chapter~4) additionally provides ongoing,
zero-user-effort second-factor-equivalent assurance throughout a
session, in contrast to MFA's point-in-time-only second factor.

\section{PKI / Smart-Card Authentication}

Public Key Infrastructure authentication binds a user identity to an
X.509 certificate, typically stored on a smart card or hardware token,
with authentication performed via a challenge-response signature using
the certificate's private key. This is the dominant model for
government and defense identity systems (e.g., India's DSC --- Digital
Signature Certificate --- regime under the IT Act 2000) and enterprise
VPN authentication. PKI signature verification (RSA-2048 or ECDSA P-256)
costs low single-digit milliseconds server-side, but the overall system
carries substantial \textit{operational} overhead: certificate issuance
requires a registration authority, certificates expire and must be
renewed, and revocation checking (via CRL or OCSP) adds a network
round-trip --- OCSP responder latency is commonly reported in the
50--200\,ms range in production CA infrastructure, comparable to
\textsc{ZK-Auth}'s full end-to-end latency for what is, in PKI's case,
only the revocation-check sub-step.

\textbf{Comparison.} \textsc{ZK-Auth}'s Groth16 verification is
$\mathcal{O}(1)$ regardless of the number of registered users or
credentials issued (Section~6.6), whereas PKI revocation checking scales
with the size of the issuer's revocation list (CRL) or requires a live
OCSP query per authentication. \textsc{ZK-Auth} additionally requires no
physical token distribution logistics, a substantial operational
simplification relative to smart-card issuance at population scale.
However, PKI's certificate-based model provides legally-recognized
non-repudiation (a signed transaction is attributable to a specific
legal identity) that \textsc{ZK-Auth}'s anonymity-preserving nullifier
design deliberately does \textit{not} provide --- this is a genuine
trade-off, not a strict improvement, and is noted as such in the
Limitations of Chapter~7.

\section{Pure OAuth 2.0 / OIDC (Without ZKP)}

It is necessary to distinguish \textsc{ZK-Auth} from a
\textit{conventional} OAuth~2.0/OpenID Connect deployment, since
Section~\ref{sec:oauth_zkp} establishes that \textsc{ZK-Auth} is itself
an OAuth Authorization Server. A conventional OAuth/OIDC identity
provider performs the
Authorization Code grant's authentication step via a
\textit{password}-backed login form --- the relying-party delegation
layer is identical to \textsc{ZK-Auth}'s, but the authentication
primitive underneath it is exactly the password mechanism critiqued in
Section~6.2. This is precisely the architectural point made in
Section~\ref{sec:oauth_zkp}: \textsc{ZK-Auth} is not an alternative to
OAuth, it is an alternative \textit{authentication primitive within}
OAuth.

A further practical concern with conventional OAuth/OIDC, noted in the
Introduction (Chapter~1), is centralization: the identity provider
learns which relying party a given user is authenticating to on every
login, and a breach of the identity provider compromises every relying
party that trusts it. The 2023 Microsoft Storm-0558 incident, in which a
compromised Microsoft signing key allowed forgery of authentication
tokens across multiple federated services, is a widely-cited recent
instance of this systemic risk class.

\textbf{Comparison.} \textsc{ZK-Auth} retains OAuth's relying-party
delegation contract (Section~\ref{sec:oauth_zkp}) while replacing the
password-backed authentication primitive with a ZKP, removing the
specific class of risk illustrated above: a breach of \textsc{ZK-Auth}'s
database exposes only Poseidon commitments (Scenario~A,
Section~\ref{sec:security}-D), not credentials usable to forge
authentication elsewhere.

\section{SAML 2.0}

Security Assertion Markup Language is the XML-based
predecessor to OAuth/OIDC in enterprise federated identity, still
widely deployed in legacy enterprise single-sign-on. SAML's
verbose XML assertion format and mandatory XML-Signature verification
impose materially higher processing overhead than OAuth's compact JSON
Web Tokens; published enterprise IAM benchmarking commonly reports SAML
assertion validation in the 100--300\,ms range due to XML canonicalization
and signature verification cost, compared to single-digit-millisecond
JWT verification. SAML shares OAuth's centralized-IdP risk profile
(Section~6.5) without OAuth's lighter-weight token format.

\textbf{Comparison.} \textsc{ZK-Auth}'s full authentication round trip
(132.51\,ms) is competitive with or faster than typical reported SAML
assertion-validation latency alone, before any network round-trip to the
IdP is even included, while providing the credential-theft resistance
SAML (like OAuth) does not natively offer.

\section{WebAuthn / FIDO2}

WebAuthn/FIDO2 is the most directly comparable modern
baseline: it eliminates shared secrets entirely, binding authentication
to a device-resident asymmetric keypair (often hardware-backed via a
secure enclave or TPM) and a biometric or PIN gesture to unlock the
private key locally. Google's large-scale passkey rollout reports
WebAuthn authentication completing in well under 200\,ms end-to-end in
production telemetry.

\textbf{Comparison.} WebAuthn and \textsc{ZK-Auth} share the core
property of never transmitting a shared secret, and report comparable
order-of-magnitude latency. The two systems differ in what they
additionally provide: WebAuthn is bound to a specific physical device's
hardware key; it offers no attribute-level selective
disclosure and no native continuous-session mechanism. \textsc{ZK-Auth}
adds both selective disclosure (Chapter~4) and behavioral continuity
(Chapter~4) on top of an equivalent no-shared-secret authentication
primitive, at the cost of being software-based rather than
hardware-key-anchored --- a genuine trade-off discussed further in
Chapter~7's Limitations, where WebAuthn-bound secret storage is
identified as a hardening direction for \textsc{ZK-Auth} itself.

\section{Behavioral-Biometric-Only Continuous Authentication}

Commercial systems such as BioCatch use
behavioral biometrics (keystroke dynamics, mouse patterns, device
interaction signals) as the \textit{primary} authentication or
fraud-detection signal, rather than as a secondary continuity check.
Commercial keystroke-biometric products report genuine-acceptance
accuracy exceeding 98\% in production deployment; published academic
keystroke-biometric systems
report comparable accuracy figures on
research benchmarks. These systems carry an inherent architectural risk
when used as a \textit{primary} authenticator: behavioral biometrics are
probabilistic and drift over time.

\textbf{Comparison.} \textsc{ZK-Auth} deliberately avoids this failure
mode by construction: the LSTM behavioral layer (Chapter~4) is
\textit{never} the sole authentication gate. It is a post-ZKP
session-continuity check; a behavioral-model failure (false negative or
false positive) degrades gracefully to ZKP-only authentication via
step-up re-proof, never to either unauthenticated access or a hard
lockout driven by behavioral drift alone. \textsc{ZK-Auth}'s measured
FAR (6.15\%, Table~\ref{tab:lstm_results}) is the security-relevant
figure precisely because false accepts at the \textit{continuity} layer
only grant continued access to a session that was already cryptographically
authenticated by ZKP --- a materially lower-stakes failure mode than a
false accept at a \textit{primary} authenticator.

\section{Consolidated Quantitative Comparison}

Table~\ref{tab:baseline_quant} consolidates the literature-reported
latency figures discussed above against \textsc{ZK-Auth}'s measured
results, and Fig.~\ref{fig:baseline_chart} plots them directly.

\begin{table}[htbp]
\renewcommand{\arraystretch}{1.25}
\caption{Consolidated Quantitative Baseline Comparison. ZK-Auth figures
are measured (this thesis, Chapter~5); all other figures are
literature-reported, sourced as cited in Sections~6.2--6.8.}
\label{tab:baseline_quant}
\centering
\footnotesize
\begin{tabular}{lcl}
\toprule
\textbf{System} & \textbf{Typical latency (ms)} & \textbf{Source} \\
\midrule
Password (Argon2id hash+verify) & 172.82 (measured) & This thesis, Table~5.3 \\
MFA (TOTP verify only, excl.\ UX) & $<$1 & RFC~6238 \\
PKI / OCSP revocation check       & 50--200 & CA/PKI industry telemetry \\
SAML 2.0 assertion validation     & 100--300 & Enterprise IAM benchmarks \\
WebAuthn/FIDO2 (full auth)        & $<$200 & FIDO Alliance / Google \\
\textbf{ZK-Auth (full E2E, Node.js)} & \textbf{132.51 (measured)} & \textbf{This thesis, Table~5.2} \\
\bottomrule
\multicolumn{3}{l}{\footnotesize MFA excludes user-perceived OTP retrieval/entry time (typically 10--20\,s).}\\
\end{tabular}
\end{table}

\begin{figure}[htbp]
\centering
\begin{tikzpicture}[x=1.9cm,y=0.018cm]
  \draw[->] (-0.4,0) -- (5.4,0) node[right,font=\scriptsize] {};
  \draw[->] (-0.4,0) -- (-0.4,330) node[above,font=\scriptsize] {ms};
  \foreach \y in {0,100,200,300}{
    \draw[gray!30] (-0.4,\y) -- (5.0,\y);
    \node[left,font=\tiny] at (-0.4,\y) {\y};
  }
  \simplebar{0}{0}{0.6}{172.82}{gray!45}{172.82}
  \simplebar{1}{0}{0.6}{125}{gray!45}{125}
  \simplebar{2}{0}{0.6}{200}{gray!45}{200}
  \simplebar{3}{0}{0.6}{180}{gray!45}{180}
  \simplebar{4}{0}{0.6}{132.51}{blue!60}{132.51}
  \node[barlabel,rotate=20,anchor=east] at (0.3,-15) {Password};
  \node[barlabel,rotate=20,anchor=east] at (1.3,-15) {PKI/OCSP};
  \node[barlabel,rotate=20,anchor=east] at (2.3,-15) {SAML};
  \node[barlabel,rotate=20,anchor=east] at (3.3,-15) {WebAuthn};
  \node[barlabel,rotate=20,anchor=east] at (4.3,-15) {ZK-Auth};
  \node[draw,fill=gray!45,minimum size=6pt] at (5.6,300) {};
  \node[barlabel,right] at (5.75,300) {Literature-reported};
  \node[draw,fill=blue!60,minimum size=6pt] at (5.6,270) {};
  \node[barlabel,right] at (5.75,270) {ZK-Auth (measured)};
\end{tikzpicture}
\caption{Consolidated latency comparison, \textsc{ZK-Auth} included
alongside all six baselines. PKI/OCSP, SAML, and WebAuthn bars use
literature-reported midpoints of the ranges in
Table~\ref{tab:baseline_quant}; Password and \textsc{ZK-Auth} bars are
measured figures from this thesis (Table~\ref{tab:zkp_results}).
\textsc{ZK-Auth}, at 132.51\,ms, is faster than every literature-reported
baseline except the trivial TOTP-verification-only figure.}
\label{fig:baseline_chart}
\end{figure}

\section{Summary of Baseline Positioning}
\label{sec:summary_baseline}

No single baseline in this chapter dominates \textsc{ZK-Auth} across
both the qualitative dimensions of Table~\ref{tab:comparison} and the
quantitative latency figures of Table~\ref{tab:baseline_quant}.
WebAuthn/FIDO2 is the closest competitor on security properties but
lacks selective disclosure and continuous authentication; PKI provides
legal non-repudiation that \textsc{ZK-Auth} deliberately forgoes in
exchange for anonymity; behavioral-biometric-only systems achieve
comparable accuracy figures but carry primary-authenticator risk that
\textsc{ZK-Auth}'s architecture avoids by design. This chapter's
contribution is not a claim that \textsc{ZK-Auth} is unconditionally
superior to every baseline on every axis, but a transparent,
side-by-side accounting of where it improves, where it matches, and
where it trades against each established alternative.

%==============================================================================
\chapter{Security Analysis, Limitations, and Future Work}
\label{sec:security}
%==============================================================================

\section{Formal Security Properties}

\begin{theorem}[Authentication Soundness]
Under the $q$-PKE assumption in the generic group model, no probabilistic
polynomial-time (PPT) adversary $\mathcal{A}$ can produce a valid Groth16
proof $\pi$ for the auth circuit without knowledge of a secret $s$ such
that $c = \mathrm{Poseidon}(s)$ matches the stored commitment, except with
probability $\mathsf{negl}(\lambda)$ for security parameter
$\lambda = 128$.
\end{theorem}

\begin{proof}
Groth16's knowledge soundness (Groth, Theorem~1) guarantees that for any
PPT prover producing an accepting proof $\pi$ for public input $n$ and
claimed outputs $(c, \eta)$, there exists a PPT extractor $\mathcal{E}$
that, given the same random tape as $\mathcal{A}$, outputs a witness $w$
satisfying the circuit's R1CS constraint system $Aw \circ Bw = Cw$. This
guarantee is generic in the constraint system: it holds for any circuit
reducible to a quadratic arithmetic program, independent of what the
constraints compute.

The auth circuit's constraint system enforces, among its 932 constraints,
two Poseidon constraint subsystems corresponding to
Eq.~\ref{eq:commitment} ($t{=}2$, binding the witness's secret component
$s$ to output $c$) and Eq.~\ref{eq:nullifier} ($t{=}3$, binding $s$ and
public input $n$ to output $\eta$). Because Poseidon's round function is
itself expressed as a sequence of R1CS constraints ($n_F{+}n_P{=}65$
rounds total per invocation, each contributing S-box and MDS-mixing
constraints), any witness $w$ extracted by $\mathcal{E}$ that satisfies
the full constraint system must in particular satisfy both Poseidon
subsystems with the \emph{same} witness value for $s$ in both. This is a
property of the circuit's design, not merely of Groth16 in the abstract:
the circuit declares a single witness signal for $s$ and routes it into
both the commitment sub-circuit and the nullifier sub-circuit, so
$\mathcal{E}$ cannot extract two different values that separately satisfy
each hash binding---there is only one variable for $\mathcal{E}$ to
extract.

It follows that extraction yields a value $s^* = w[\mathtt{secret}]$
satisfying $\mathrm{Poseidon}(s^*) = c$ exactly when $c$ is the proof's
claimed output, i.e., the same $s^*$ an honest prover would need to know
to generate $\pi$. If $c$ matches the value stored at registration, then
by Poseidon's collision and preimage resistance at 128-bit security, $s^*
= s$ with overwhelming probability: any $s^* \neq s$ satisfying
$\mathrm{Poseidon}(s^*) = \mathrm{Poseidon}(s)$ would itself constitute a
Poseidon collision. Thus an adversary producing an accepting $\pi$ without
knowledge of $s$ implies, via $\mathcal{E}$, either a break of the
$q$-PKE assumption (extraction itself fails) or a Poseidon collision---
both events of probability $\mathsf{negl}(\lambda)$ for $\lambda=128$.
\end{proof}

\begin{theorem}[Zero-Knowledge]
Groth16 is computationally zero-knowledge: any proof $\pi$ produced by an
honest prover is computationally indistinguishable from a simulator's
output produced without knowledge of $s$.
\end{theorem}

\begin{proof}[Proof sketch]
Follows directly from Groth's Theorem~3, which exhibits a simulator that,
given only the public inputs and the (public) proving/verification key
pair, produces a proof distribution computationally indistinguishable
from that of an honest prover holding the witness. Applied to the auth
circuit, unlinkability across sessions follows as a corollary: because
distinct challenge nonces $n$ produce distinct nullifiers
$\eta = \mathrm{Poseidon}(s, n)$ for the same fixed secret $s$, and
because Poseidon is modeled as a random oracle for this purpose, an
observer with no knowledge of $s$ cannot link nullifiers $\eta_1$ and
$\eta_2$ arising from two different authentication sessions of the same
user, since doing so would require either recovering $s$ from $\eta_1$
(a Poseidon preimage break) or distinguishing the simulated proof
distribution from the honest one (a break of the zero-knowledge property
itself).
\end{proof}

\begin{corollary}[Nullifier Replay Prevention]
A previously submitted, valid nullifier $\eta$ cannot be accepted a
second time by the verifier.
\end{corollary}

\begin{proof}[Justification]
This follows not from the zero-knowledge proof system itself but from the
server-side protocol construction: nullifiers are atomically inserted
into the Redis nullifier set via \texttt{SADD} under a distributed lock
\emph{before} token issuance, with a \texttt{SISMEMBER} pre-check
rejecting any nullifier already present in the set. Because this check
occurs prior to the (comparatively expensive) Groth16 verification step,
a replayed $(\pi, \eta, c)$ tuple is rejected at negligible computational
cost, independent of whether $\pi$ would otherwise verify correctly.
Combined with the unlinkability property of Theorem~2 (an attacker cannot
predict a future session's nullifier from a captured one), this closes
the replay attack surface entirely for the captured-and-resubmitted proof
scenario.
\end{proof}

\section{Attack Scenario Walkthrough}

To make the abstract threat model of Chapter~3 concrete, we trace three
representative attack attempts against the deployed system and the
corresponding defense each scenario exercises.

\textbf{Scenario A --- Database exfiltration.} An attacker obtains a full
dump of the PostgreSQL \texttt{users} table. The table contains only
Poseidon commitments $c$, never the secret $s$. Per Theorem~1, recovering
$s$ from $c$ requires inverting Poseidon, computationally infeasible at
128-bit security. The attacker gains nothing usable for authentication,
in sharp contrast to a conventional bcrypt/Argon2id database breach,
which at minimum exposes hashes vulnerable to offline dictionary attack
against weak user-chosen passwords.

\textbf{Scenario B --- Proof replay.} An attacker intercepts a valid
$(\pi, \eta, c)$ tuple in transit and attempts to resubmit it later. Per
Corollary~1, the nullifier $\eta$ is already present in the Redis
nullifier set from the original submission; the \texttt{SISMEMBER}
pre-check rejects the replay before proof verification is even
attempted, at negligible computational cost to the server.

\textbf{Scenario C --- Session hijacking post-authentication.} An
attacker steals a valid JWT (e.g., via cross-site scripting) and begins
issuing requests as the victim. The behavioral telemetry stream for the
hijacked session deviates from the victim's learned profile; once the
EMA-smoothed risk score crosses 0.90, the server forces step-up
re-authentication (Chapter~3), which the attacker cannot complete without
the victim's registration secret $s$, terminating the hijacked session
within the detection window characterized in Chapter~5.

\section{Resolution of Research Challenges}

\textbf{Challenge 1 (Proof complexity and cross-environment viability):}
Addressed through Circom 2.x circuit optimization (932 constraints) and
Web Worker / WebView WASM execution, with empirical validation across
four environments (Chapter~5) demonstrating sub-second proving latency
even on the slowest measured hardware tier.

\textbf{Challenge 2 (Attribute privacy):} Resolved through Poseidon
Merkle selective disclosure (Chapter~4); verifiers receive only boolean
predicate results, never attribute values, with the depth-8 design choice
justified quantitatively in Section~5.7.

\textbf{Challenge 3 (Session continuity):} Addressed by the LSTM
behavioral biometric layer with ZKP step-up re-authentication
(Chapter~4), trained and evaluated on authentic (not synthetic)
behavioral data, achieving AUC-ROC 0.8157 (Chapter~5).

\textbf{Challenge 4 (Empirical and theoretical rigor):} Addressed by the
four-environment benchmark suite (Chapter~5), the ablation studies
justifying the Merkle depth and LSTM window hyperparameters
(Section~5.7), and the extraction-based soundness reduction (Theorem~1)
tying the specific Poseidon-bound circuit constraints to Groth16's
generic-group guarantees, rather than relying on a citation alone.

\section{Limitations}

\textit{(i) Single-feature behavioral coverage:} The Balabit dataset
provides mouse dynamics only; the live telemetry pipeline's six-feature
vector includes keyboard dwell and scroll delta, which are zero
throughout the current training corpus. Extending training to a combined
dataset (e.g., Balabit plus a keystroke-dynamics corpus) is identified as
near-term future work.

\textit{(ii) Trusted setup:} the public Hermez transcript is used rather
than a deployment-specific ceremony; while broad public participation in
the original ceremony provides strong practical assurance, a dedicated
ceremony would further strengthen the trust assumption underlying
Theorem~1's $q$-PKE premise.

\textit{(iii) Secret storage:} browser \texttt{localStorage} is
accessible to JavaScript on the same origin; production deployment should
migrate to WebAuthn-bound or hardware security key storage for the
registration secret.

\textit{(iv) FRR/FAR tradeoff:} the 0.75 operating threshold is tuned
conservatively toward minimizing false acceptance at the cost of elevated
false rejection (Table~\ref{tab:lstm_results}); this is a deliberate
security-first choice, but the resulting re-authentication friction for
legitimate users with atypical mouse behavior is not separately quantified
as a UX metric in this thesis.

\textit{(v) Ablation depth:} the Merkle-depth and LSTM-window ablations
in Section~5.7 use measured depth-8/window-50 figures with linearly
derived estimates at the alternative settings; a fully rigorous ablation
would recompile the disclosure circuit at depth-4/16 and retrain the LSTM
at window-20/100 to obtain measured (rather than derived) figures at
every setting, which is identified as the most immediately actionable
extension of this evaluation.

\textit{(vi) Multi-institution load:} the issuance layer (Chapter~4) is
functionally validated for single-institution onboarding and credential
signing but has not been load-tested under simultaneous multi-tenant
traffic.

\section{Future Directions}

Future extensions to this work include: completing the Merkle-depth and
LSTM-window ablations with measured (rather than derived) data at all
settings; multi-modal behavioral training incorporating keyboard and
scroll telemetry alongside mouse dynamics; iOS WKWebView benchmarking to
complete cross-platform mobile coverage; lattice-based (LWE) post-quantum
ZKP constructions as a long-term hedge against quantum cryptanalysis of
BN254's discrete-logarithm assumption; FPGA-based multi-scalar-
multiplication acceleration for edge-device proving, building on the
acceleration techniques surveyed in Chapter~2~\cite{butt2024}; threshold
multi-party signing for institutional key custody, removing the single-
point-of-trust inherent in a single Ed25519 signing key per institution;
and integration of \textsc{ZK-Auth} authentication with a complementary
attribute-based encryption access-control layer, as motivated by the
ABE survey in Chapter~2~\cite{prasad2025abereview}.

%==============================================================================
\chapter{Conclusion}
%==============================================================================

This thesis presented \textsc{ZK-Auth}, a deployable authentication
infrastructure combining a 932-constraint Groth16 zero-knowledge proof
circuit over BN254, Poseidon Merkle selective disclosure, LSTM behavioral
biometric session continuity trained on authentic public telemetry data,
and multi-tenant Ed25519 institutional credential issuance. The system
was implemented end-to-end and empirically evaluated across four
heterogeneous deployment environments---a Node.js server, a Chrome
browser client, and two Android devices spanning a flagship-to-midrange
hardware spread---demonstrating a median end-to-end authentication
latency of 132.51\,ms on the server reference environment (23\% faster
than an Argon2id password-hashing baseline) and proof-generation times
remaining within sub-second bounds across the full measured client
hardware range, from 61.93\,ms (Node.js) to 361.0\,ms (mid-tier Android
tablet).

The behavioral biometric subsystem, evaluated on the public Balabit
Mouse Dynamics Challenge dataset rather than synthetic data, achieved an
AUC-ROC of 0.8157 with a 6.15\% false-acceptance rate at the operating
threshold, and the ablation studies of Chapter~5 provided a quantitative
rather than purely qualitative justification for the system's central
architectural hyperparameters---Merkle tree depth and LSTM window length.
A formal extraction-based reduction (Theorem~1) demonstrated explicitly
how the auth circuit's shared-witness construction preserves Groth16's
generic-group soundness guarantee for the specific Poseidon-bound
commitment and nullifier bindings used in this system, going beyond a
citation-only security argument.

Taken together, these results demonstrate that zero-knowledge-proof-based
authentication is not merely privacy-preserving in theoretical
construction but practically deployable at production-acceptable latency
across the realistic spread of consumer client hardware, with direct
implications for cloud security architecture, decentralized identity
infrastructure, and regulatory compliance under India's DPDP Act 2023 and
the EU's GDPR. The complete implementation, including circuit sources,
compiled proving artifacts, benchmark scripts, and the Flutter mobile
client, is maintained in the project's public repository to support
independent reproduction of every quantitative claim made in this thesis.

%==============================================================================
% PUBLICATIONS
%==============================================================================
\clearpage
\chapter*{Publications}
\addcontentsline{toc}{chapter}{Publications}

\begin{enumerate}[label={[\arabic*]}, itemsep=1.5ex, leftmargin=*]
  \item Abhay Dandge, Sameeksha Prasad, and Namita Tiwari. ``A Review on Zero Knowledge Proof For Authentication In Cloud Computing.'' \textit{FrontSci-2025, The National Conference on Frontier Research in Sciences}, Maulana Azad National Institute of Technology. \textbf{(Presented)}

  \item Abhay Dandge, Sameeksha Prasad, and Namita Tiwari. ``Zero Knowledge Proof For Authentication In Cloud Computing : A Review.'' \textit{10th International Conference on Internet of Things and Connected Technologies (ICIoTCT) 2025}. \textbf{(Accepted)}

  \item Abhay Dandge, Sameeksha Prasad, and Namita Tiwari. ``A Comprehensive Review of Zero-Knowledge Proofs in Attribute-Based Encryption Frameworks.'' \textit{First International Conference on Nexus of Digitalization, Intelligence and Applications 2026}. \textbf{(Accepted)}

  \item Abhay Dandge, Sameeksha Prasad, and Namita Tiwari. ``ZK-Auth: A Zero-Knowledge Proof-Based Authentication Framework with Selective Disclosure and Continuous Behavioral Biometric Verification.''
\end{enumerate}

%==============================================================================
% BIBLIOGRAPHY
%==============================================================================
\clearpage
\addcontentsline{toc}{chapter}{Bibliography}
\bibliographystyle{unsrt}
\begin{thebibliography}{99}

\bibitem{goldreich94}
O. Goldreich and Y. Oren, ``Definitions and properties of zero-knowledge proof
systems,'' \textit{J. Cryptology}, vol.~7, no.~1, pp.~1--32, 1994.

\bibitem{bfm88}
M. Blum, P. Feldman, and S. Micali, ``Non-interactive zero-knowledge and its
applications,'' in \textit{Proc. 20th ACM STOC}, 1988, pp.~103--112.

\bibitem{fiatshamir}
A. Fiat and A. Shamir, ``How to prove yourself: Practical solutions to identification
and signature problems,'' in \textit{Proc. CRYPTO 1986}, LNCS vol.~263, 1987,
pp.~186--194.

\bibitem{parno13}
B. Parno, J. Howell, C. Gentry, and M. Raykova, ``Pinocchio: Nearly practical
verifiable computation,'' in \textit{Proc. IEEE S\&P}, 2013, pp.~238--252.

\bibitem{gennaro13}
R. Gennaro, C. Gentry, B. Parno, and M. Raykova, ``Quadratic span programs and
succinct NIZKs without PCPs,'' in \textit{Proc. EUROCRYPT 2013}, LNCS vol.~7881,
pp.~626--645.

\bibitem{groth16}
J. Groth, ``On the size of pairing-based non-interactive arguments,'' in
\textit{Advances in Cryptology -- EUROCRYPT 2016}, LNCS vol.~9666, Springer,
pp.~305--326, 2016.

\bibitem{zcash14}
E. Ben-Sasson et al., ``Zerocash: Decentralized anonymous payments from Bitcoin,''
in \textit{Proc. IEEE S\&P}, 2014, pp.~459--474.

\bibitem{plonk}
A. Gabizon, Z. J. Williamson, and O. Ciobanu, ``PLONK: Permutations over
Lagrange-bases for Oecumenical Noninteractive arguments of Knowledge,''
\textit{IACR ePrint} 2019/953, 2019.

\bibitem{starks}
E. Ben-Sasson et al., ``Scalable, transparent, and post-quantum secure computational
integrity,'' \textit{IACR ePrint} 2018/046, 2018.

\bibitem{bulletproofs}
B. B\"unz, J. Bootle, D. Boneh, A. Poelstra, P. Wuille, and G. Maxwell,
``Bulletproofs: Short proofs for confidential transactions and more,'' in
\textit{Proc. IEEE S\&P}, 2018, pp.~315--334.

\bibitem{dandge2025zkpcloud}
A. Dandge, S. Prasad, N. Tiwari, and M. Chawla, ``Zero-Knowledge Proof for
Authentication in Cloud Computing: A Review,'' in \textit{Proc. FrontSci 2025},
Springer Nature, 2025.

\bibitem{khandait2024}
A. U. Khandait and T. R. Pattanshetti, ``Lightweight authentication system based
on ZKP for IoT devices,'' in \textit{Proc. GCCIT 2024}, IEEE, 2024.

\bibitem{pathak2021}
A. Pathak et al., ``Secure authentication using zero knowledge proof,''
in \textit{Proc. ASIANCON 2021}, IEEE, 2021.

\bibitem{butt2024}
S. A. Butt et al., ``if-ZKP: Intel FPGA-based acceleration of zero knowledge
proofs,'' \textit{arXiv:2412.12481}, 2024.

\bibitem{hou2022}
D. Hou et al., ``Privacy-preserving energy trading using blockchain and ZKP,''
in \textit{Proc. IEEE Blockchain}, 2022.

\bibitem{verma2024}
P. Verma, V. Tripathi, and B. Pant, ``ZeroMedChain: ZKP integration for healthcare
identity and access management,'' in \textit{Proc. INDIACom}, IEEE, 2024.

\bibitem{naidu2023}
D. Naidu et al., ``Smart contract for privacy preserving authentication using ZKP,''
in \textit{Proc. OCIT}, IEEE, 2023.

\bibitem{sahai05}
A. Sahai and B. Waters, ``Fuzzy identity-based encryption,'' in
\textit{Proc. EUROCRYPT 2005}, LNCS vol.~3494, Springer, pp.~457--473, 2005.

\bibitem{goyal06}
V. Goyal, O. Pandey, A. Sahai, and B. Waters, ``Attribute-based encryption for
fine-grained access control of encrypted data,'' in \textit{Proc. 13th ACM CCS},
2006, pp.~89--98.

\bibitem{kate10}
A. Kate, G. M. Zaverucha, and I. Goldberg, ``Constant-size commitments to
polynomials and their applications,'' in \textit{Proc. ASIACRYPT 2010},
pp.~177--194.

\bibitem{prasad2025abereview}
S. Prasad, A. Dandge, N. Tiwari, and M. Chawla, ``A Comprehensive Review of
Zero-Knowledge Proofs in Attribute-Based Encryption Frameworks,'' accepted,
\textit{Proc. ICNDIA 2026}, IEEE, 2026.

\bibitem{yang2024}
R. Yang, Y. Liu, and S. Chen, ``Advancing user privacy in virtual power plants:
A novel ZKP-based distributed attribute encryption approach,''
\textit{Electronics}, vol.~13, no.~7, pp.~1283--1298, 2024.

\bibitem{ren2024}
Z. Ren, H. Li, and J. Wu, ``Blockchain-based CP-ABE data sharing using distributed
KMS and ZKP,'' \textit{J. King Saud Univ. Comput. Inf. Sci.}, vol.~36, no.~3, 2024.

\bibitem{bbs}
D. Boneh, X. Boyen, and H. Shacham, ``Short group signatures,'' in
\textit{Proc. CRYPTO 2004}, LNCS vol.~3152, Springer, pp.~41--55, 2004.

\bibitem{sdjwt}
O. Terbu, D. Fett, and B. Campbell, ``Selective Disclosure for JWTs (SD-JWT),''
IETF Internet-Draft, 2024.

\bibitem{cl}
J. Camenisch and A. Lysyanskaya, ``An efficient system for non-transferable anonymous
credentials with optional anonymity revocation,'' in \textit{Proc. EUROCRYPT 2001},
pp.~93--118, 2001.

\bibitem{shepherd95}
S. Shepherd, ``Continuous authentication by analysis of keyboard typing
characteristics,'' in \textit{Proc. ECSD}, 1995, pp.~111--114.

\bibitem{nakkabi11}
A. Nakkabi, I. Trabelsi, and A. Chikh, ``Improving mouse dynamics biometric
performance using variance reduction,'' \textit{IEEE Trans. SMC}, vol.~41, no.~1,
pp.~42--48, 2011.

\bibitem{frank13}
M. Frank et al., ``Touchalytics: On the applicability of touchscreen input as a
behavioral biometric,'' \textit{IEEE Trans. IFS}, vol.~8, no.~1, pp.~136--148, 2013.

\bibitem{typenet}
P. Acien et al., ``TypeNet: Deep learning keystroke biometrics,''
\textit{IEEE Trans. BTAS}, vol.~4, no.~1, pp.~57--67, 2022.

\bibitem{polisetti2024}
V. K. Polisetti et al., ``zkSHIELD: A web-based cross-chain zero knowledge asset
proof,'' in \textit{Proc. ICCCNT}, IEEE, 2024.

\bibitem{poseidon}
L. Grassi et al., ``Poseidon: A new hash function for zero-knowledge proof systems,''
in \textit{Proc. USENIX Security}, 2021, pp.~519--535.

\bibitem{bernstein2011}
D. J. Bernstein et al., ``High-speed high-security signatures,'' in \textit{Proc.
CHES 2011}, LNCS vol.~6917, Springer, pp.~124--142, 2011.

\bibitem{hochreiter97}
S. Hochreiter and J. Schmidhuber, ``Long short-term memory,''
\textit{Neural Computation}, vol.~9, no.~8, pp.~1735--1780, 1997.

\bibitem{argon2}
A. Biryukov, D. Dinu, and D. Khovratovich, ``Argon2: New generation of memory-hard
functions for password hashing and other applications,'' in \textit{Proc. IEEE
EuroS\&P}, 2016, pp.~292--302.

\bibitem{saml}
S. Cantor, J. Kemp, R. Philpott, and E. Maler, ``Assertions and Protocols for
the OASIS Security Assertion Markup Language (SAML) V2.0,'' OASIS Standard, 2005.

\end{thebibliography}

\end{document}